\documentclass[8pt, a4paper]{article}

% Page Geometry: Compact margins with proper space reserved for tiny header
\usepackage[landscape, margin=0.35in, top=0.45in, bottom=0.35in, headheight=9pt, headsep=6pt]{geometry}

% Multi-column Layout
\usepackage{multicol}
\setlength{\columnsep}{18pt}       % gap between columns
\setlength{\columnseprule}{0.4pt}  % Column divider line

% Code Formatting Setup
\usepackage{listings}
\usepackage{xcolor}
\usepackage{textcomp}              % straight quotes ' in code (upquote)

\definecolor{codegray}{rgb}{0.5,0.5,0.5}
\definecolor{codepurple}{rgb}{0.58,0.0,0.82}

\lstset{
    aboveskip=0pt,
    belowskip=0pt,
    lineskip=0pt,
    language=C++,
    basicstyle=\ttfamily\footnotesize,  % same font size as before
    columns=fixed,
    basewidth=4.25pt,                   % natural width of the 8pt typewriter font: 57 characters per line
    keepspaces=true,
    upquote=true,
    keywordstyle=\color{blue}\bfseries,
    commentstyle=\color{codegray}\itshape,
    stringstyle=\color{codepurple},
    numberstyle=\tiny\color{codegray},
    numbers=left,
    stepnumber=1,
    numbersep=3pt,
    xleftmargin=9pt,                    % room for the line numbers inside the column
    showstringspaces=false,
    breaklines=true,                    % safety net only: every code line is at most 57 characters
    breakatwhitespace=false,
    tabsize=2
}

% Micro-spacing adjustment for section headings
\usepackage{titlesec}
\titlespacing*{\section}{0pt}{1pt}{1pt}
\titlespacing*{\subsection}{0pt}{1pt}{0pt}

% Header/Footer Setup (Ultra-compact header)
\usepackage{fancyhdr}
\pagestyle{fancy}
\fancyhf{}
\lhead{\fontsize{7}{8}\selectfont JU\_Zer0ne \ \textbar\ Jahangirnagar University}
\rhead{\fontsize{7}{8}\selectfont Page \thepage}
\renewcommand{\headrulewidth}{0.2pt}

% Table of Contents formatting for 3 columns
\usepackage{tocloft}
\renewcommand{\cftsecfont}{\footnotesize\bfseries}
\renewcommand{\cftsecpagefont}{\footnotesize\bfseries}
\renewcommand{\cftsubsecfont}{\scriptsize}
\renewcommand{\cftsubsecpagefont}{\scriptsize}

% Tight spacing to maximize vertical space per column
\setlength{\cftbeforesecskip}{0pt}
\setlength{\cftbeforesubsecskip}{-2pt}

% Minimal indentation to maximize horizontal text area in 3 columns
\setlength{\cftsecindent}{0pt}
\setlength{\cftsubsecindent}{4pt}

\begin{document}

% =========================================================
% PAGE 1: COVER & AUTOMATIC INDEX (3 COLUMNS)
% =========================================================
\thispagestyle{empty}

\begin{center}
    {\Large \bfseries JU\_Zer0ne}\\[3pt]
    {\small Jahangirnagar University}\\[3pt]
    {\textbf{Members:} Tanvir Alam Shihab, Member 2, Member 3}
\end{center}

\vspace{-3pt}
\hrule
\vspace{6pt}

\begin{center}
    {\small \bfseries Algorithm Index}
\end{center}

\vspace{-5pt}

% 3-Column Index Layout
\begin{multicols*}{3}
\addtocontents{toc}{\vspace{-5pt}}
\addtocontents{toc}{\vspace{-5pt}}
\addtocontents{toc}{\vspace{-5pt}}
\tableofcontents
\end{multicols*}

\newpage
\setcounter{page}{2}

% =========================================================
% SUBSEQUENT PAGES: ALGORITHMS CHEATSHEET
% Every snippet assumes the Template in section 1
% (#define int long long, all(x), MOD, inf, rng).
% =========================================================
\begin{multicols*}{3}
\section{General}
\subsection{Template}
\begin{lstlisting}
#include <bits/stdc++.h>
#include <ext/pb_ds/assoc_container.hpp>
#include <ext/pb_ds/tree_policy.hpp>
using namespace std;
using namespace __gnu_pbds;
typedef long long ll;
typedef unsigned long long ull;
// NOTE: every snippet in this notebook assumes this:
// int is 64-bit, so write "signed main" and 1LL << x.
#define int ll
#define endl "\n"
#define all(x) (x).begin(), (x).end()
#define rall(x) (x).rbegin(), (x).rend()
#define YES cout << "YES\n"
#define NO cout << "NO\n"
const int MOD = 1e9 + 7;
const int inf = 1e18;
// Ordered set (PBDS):
//  s.find_by_order(k) -> iterator to k-th (0-based)
//  s.order_of_key(x)  -> how many elements are < x
//  duplicates: store pair<int,int> {value, unique id}
template <class T>
using ordered_set = tree<T, null_type, less<T>,
  rb_tree_tag, tree_order_statistics_node_update>;
// Random: rng() % n, shuffle(all(v), rng)
mt19937_64 rng(
  chrono::steady_clock::now().time_since_epoch()
    .count());

void solve() {
}
signed main() {
  ios_base::sync_with_stdio(false);
  cin.tie(NULL);
  int tc = 1;
  cin >> tc; // remove this line for a single test
  for (int cs = 1; cs <= tc; cs++) {
    // cout << "Case " << cs << ": ";
    solve();
  }
  return 0;
}
// Build:
//  g++ -std=c++17 -O2 -Wall a.cpp -o a && ./a < in.txt
// Run time (print to cerr, judge ignores it):
//  auto t0 = chrono::high_resolution_clock::now();
//  ... code ...
//  auto t1 = chrono::high_resolution_clock::now();
//  cerr << chrono::duration<double>(t1 - t0).count();
\end{lstlisting}
\subsection{Comparator for priority queue}
\begin{lstlisting}
// priority_queue: cmp(a, b) == true means a goes BELOW b
// (opposite of sort). top() is the "largest" element.
struct Compare {
  bool operator()(const vector<int> &a,
                  const vector<int> &b) const {
    // bigger a[0] comes out first
    if (a[0] != b[0]) return a[0] < b[0];
    return a[1] > b[1];        // tie: small a[1] on top
  }
};
// priority_queue<vector<int>, vector<vector<int>>,
//                Compare> pq;
// min-heap: priority_queue<T, vector<T>, greater<T>> pq;
\end{lstlisting}
\subsection{Next and previous greater element}
\begin{lstlisting}
// O(n). res[i] = index of the nearest greater element on
// the right (n if none). For "smaller" flip the sign >.
vector<int> nextGreater(const vector<int> &a) {
  int n = a.size();
  vector<int> res(n, n);
  stack<int> s; // indices still waiting for an answer
  for (int i = 0; i < n; ++i) {
    while (!s.empty() && a[i] > a[s.top()]) {
      res[s.top()] = i;
      s.pop();
    }
    s.push(i);
  }
  return res;
}
// res[i] = nearest greater on the left (-1 if none)
vector<int> prevGreater(const vector<int> &a) {
  int n = a.size();
  vector<int> res(n, -1);
  stack<int> s;
  for (int i = 0; i < n; ++i) {
    while (!s.empty() && a[s.top()] <= a[i]) s.pop();
    if (!s.empty()) res[i] = s.top();
    s.push(i);
  }
  return res;
}
// Largest rectangle under a histogram. O(n)
int largestRectangle(vector<int> h) {
  h.push_back(0); // closes every open bar at the end
  stack<int> s;   // indices with increasing heights
  int best = 0;
  for (int i = 0; i < (int)h.size(); i++) {
    while (!s.empty() && h[s.top()] >= h[i]) {
      int height = h[s.top()];
      s.pop();
      int left = s.empty() ? -1 : s.top();
      best = max(best, height * (i - left - 1));
    }
    s.push(i);
  }
  return best;
}
\end{lstlisting}
\subsection{2D spiral}
\begin{lstlisting}
struct Spiral {
  // direction order: right, down, left, up
  int dx[4] = {0, 1, 0, -1}, dy[4] = {1, 0, -1, 0};
  // n x m grid, 1..n*m from (0,0) clockwise going inward
  vector<vector<int>> in(int n, int m) {
    vector<vector<int>> r(n, vector<int>(m));
    int x = 0, y = 0, d = 0;
    for (int i = 1; i <= n * m; ++i) {
      r[x][y] = i;
      int nx = x + dx[d], ny = y + dy[d];
      if (nx < 0 || nx >= n || ny < 0 || ny >= m ||
          r[nx][ny])
        d = (d + 1) % 4; // wall or filled cell: turn
      x += dx[d], y += dy[d];
    }
    return r;
  }
  // n x n grid, 1..n*n from the center going outward
  vector<vector<int>> out(int n) {
    vector<vector<int>> r(n, vector<int>(n));
    int x = n / 2, y = n / 2, d = 0, s = 1, v = 1;
    r[x][y] = v++;
    while (v <= n * n) {
      // step lengths go 1,1,2,2,3,3,...
      for (int i = 0; i < 2; ++i, d = (d + 1) % 4)
        for (int j = 0; j < s; ++j) {
          x += dx[d], y += dy[d];
          if (x >= 0 && x < n && y >= 0 && y < n)
            r[x][y] = v++;
        }
      s++;
    }
    return r;
  }
};
\end{lstlisting}
\subsection{Counting pair sum (FFT)}
\begin{lstlisting}
// cnt[s] = number of pairs i < j with a[i] + a[j] == s.
// Values must be >= 0. O(M log M), M = max value.
// Needs multiply() from FFT.
vector<ll> countPairSum(const vector<int> &a) {
  int mx = *max_element(all(a));
  vector<int> freq(mx + 1, 0);
  for (int x : a) freq[x]++;
  vector<ll> cnt = multiply(freq, freq);
  for (int i = 0; i <= mx; i++)
    cnt[2 * i] -= freq[i];    // remove pairs with i == j
  for (auto &x : cnt) x /= 2; // (i,j) and (j,i) are same
  return cnt;
}
\end{lstlisting}
\subsection{Sprague Grundy}
\begin{lstlisting}
// Game = independent piles. grundy(state) = mex of the
// grundy values of all states reachable in ONE move.
// XOR of grundy of all piles != 0 <=> first player wins.
// Plain Nim: grundy(pile) = pile size.
const int MAXG = 1005;
int memo[MAXG]; // memset(memo, -1, sizeof(memo)) first
int grundy(int n) {
  if (n == 0) return 0; // no move left: losing state
  if (memo[n] != -1) return memo[n];
  set<int> s;
  // CHANGE: insert grundy of every reachable state.
  // Example game: remove 1, 2 or 3 stones.
  for (int t = 1; t <= 3 && t <= n; t++)
    s.insert(grundy(n - t));
  int mex = 0;
  while (s.count(mex)) mex++;
  return memo[n] = mex;
}
\end{lstlisting}
\subsection{Stars and bars}
\begin{lstlisting}
// x1 + ... + xk = sum, xi >= 0 : nCr(sum + k - 1, k - 1)
// x1 + ... + xk = sum, xi >= 1 : nCr(sum - 1, k - 1)
// With upper bound 0 <= xi <= up (inclusion-exclusion):
// Needs nCr().
int starBarUpper(int sum, int k, int up) {
  int ans = 0;
  for (int i = 0; i <= k; i++) { // i variables exceed up
    int rem = sum - (up + 1) * i;
    if (rem < 0) break;
    int cur = nCr(k, i) * nCr(rem + k - 1, k - 1) % MOD;
    if (i & 1) ans = (ans - cur + MOD) % MOD;
    else ans = (ans + cur) % MOD;
  }
  return ans;
}
\end{lstlisting}
\subsection{Binary search and ternary search}
\begin{lstlisting}
// ok() looks like F F F T T T. Returns the FIRST x in
// [lo, hi] with ok(x) true, or hi + 1 if there is none.
int firstTrue(int lo, int hi, function<bool(int)> ok) {
  hi++;
  while (lo < hi) {
    int mid = lo + (hi - lo) / 2;
    if (ok(mid)) hi = mid;
    else lo = mid + 1;
  }
  return lo;
}
// ok() looks like T T T F F F. Returns the LAST x with
// ok(x) true, or lo - 1 if there is none.
int lastTrue(int lo, int hi, function<bool(int)> ok) {
  lo--;
  while (lo < hi) {
    int mid = lo + (hi - lo + 1) / 2;
    if (ok(mid)) lo = mid;
    else hi = mid - 1;
  }
  return lo;
}
// Real numbers: do a fixed number of rounds instead.
//  for (int it = 0; it < 100; it++) {
//    double mid = (lo + hi) / 2;
//    if (ok(mid)) hi = mid; else lo = mid;
//  }
// Ternary search: position of the minimum of f on
// [lo, hi] when f strictly goes down and then up.
// For a maximum compare with >= instead of <=.
int ternaryMin(int lo, int hi, function<int(int)> f) {
  while (hi - lo > 2) {
    int m1 = lo + (hi - lo) / 3, m2 = hi - (hi - lo) / 3;
    if (f(m1) <= f(m2)) hi = m2;
    else lo = m1;
  }
  int best = lo;
  for (int x = lo + 1; x <= hi; x++)
    if (f(x) < f(best)) best = x;
  return best;
}
\end{lstlisting}
\subsection{Custom hash, compression, sliding window}
\begin{lstlisting}
// Safe hash: unordered_map can not be hacked into O(n).
struct custom_hash {
  static ull splitmix64(ull x) {
    x += 0x9e3779b97f4a7c15;
    x = (x ^ (x >> 30)) * 0xbf58476d1ce4e5b9;
    x = (x ^ (x >> 27)) * 0x94d049bb133111eb;
    return x ^ (x >> 31);
  }
  size_t operator()(ull x) const {
    static const ull FIXED =
      chrono::steady_clock::now().time_since_epoch()
        .count();
    return splitmix64(x + FIXED);
  }
};
// unordered_map<int, int, custom_hash> mp;
// gp_hash_table<int, int, custom_hash> mp; (faster)

// Coordinate compression: a[i] becomes its rank 0..k-1.
// Returns the sorted different values (rank -> value).
vector<int> compress(vector<int> &a) {
  vector<int> v = a;
  sort(all(v));
  v.erase(unique(all(v)), v.end());
  for (int &x : a)
    x = lower_bound(all(v), x) - v.begin();
  return v;
}
// Minimum of every window of size k. O(n).
// res[i] = min(a[i..i+k-1]). For maximum flip >= to <=.
vector<int> slidingMin(const vector<int> &a, int k) {
  deque<int> dq; // indices, their values are increasing
  vector<int> res;
  for (int i = 0; i < (int)a.size(); i++) {
    while (!dq.empty() && a[dq.back()] >= a[i])
      dq.pop_back();
    dq.push_back(i);
    if (dq.front() <= i - k) dq.pop_front(); // too old
    if (i >= k - 1) res.push_back(a[dq.front()]);
  }
  return res;
}
// print a __int128 value
void print128(__int128 x) {
  if (x < 0) cout << '-', x = -x;
  if (x > 9) print128(x / 10);
  cout << (char)('0' + (int)(x % 10));
}
\end{lstlisting}
\subsection{Stress testing and checklist}
\begin{lstlisting}[language={}]
# a.cpp = your solution, brute.cpp = slow but surely
# correct, gen.cpp = prints one small random test and
# takes the seed from argv[1]:
#   mt19937 rng(atoi(argv[1]));
g++ -O2 a.cpp -o a
g++ -O2 brute.cpp -o brute
g++ -O2 gen.cpp -o gen
for i in $(seq 1 100000); do
  ./gen $i > in.txt
  ./a < in.txt > out1.txt
  ./brute < in.txt > out2.txt
  if ! cmp -s out1.txt out2.txt; then
    echo "DIFFERENT on test $i"; break
  fi
done

Before submitting / when you get Wrong Answer:
 - overflow? (int is ll here, but check products > 1e18)
 - array size, N + 5, 0-based vs 1-based
 - reset globals between test cases (not with memset
   of the full array when sum of n is bounded: TLE)
 - n = 1, empty input, all equal, max values
 - printed "Case x:" format, YES/NO spelling, endl
 - a % MOD can be negative: (a % MOD + MOD) % MOD
 - double: print with setprecision(10), avoid == on it
 - reading until EOF: while (cin >> n)
 - TLE: endl -> "\n", pass vectors by reference,
   unordered_map -> custom_hash or sort + binary search
\end{lstlisting}
\section{Bit Manipulation}
\subsection{Bit tricks}
\begin{lstlisting}
// int is 64-bit here, so always shift 1LL (not 1).
bool isSet(int n, int x) { return (n >> x) & 1; }
void setBit(int &n, int x) { n |= (1LL << x); }
void flipBit(int &n, int x) { n ^= (1LL << x); }
void clearBit(int &n, int x) { n &= ~(1LL << x); }
bool isPowerOfTwo(int n) { return n && !(n & (n - 1)); }
// is n divisible by 2^k ?
bool divByPow2(int n, int k) {
  return (n & ((1LL << k) - 1)) == 0;
}
// n & (n - 1) : remove the lowest set bit
// n & (-n)    : keep only the lowest set bit
// n & (n + 1) : clear all trailing ones
// n | (n + 1) : set the lowest zero bit
// __builtin_popcountll(n) : number of set bits
// __builtin_clzll(n) : leading zeros (n > 0)
//    63 - __builtin_clzll(n) = index of highest set bit
// __builtin_ctzll(n) : index of lowest set bit (n > 0)
// all submasks of m: for (s = m; s; s = (s - 1) & m)
// a + b = (a ^ b) + 2 * (a & b) = (a | b) + (a & b)

// total number of set bits in 1, 2, ..., n
int countSetBitsUpto(int n) {
  int cnt = 0;
  while (n > 0) {
    int x = 63 - __builtin_clzll(n); // highest set bit
    if (x > 0) cnt += x * (1LL << (x - 1)); // below 2^x
    n -= 1LL << x;
    cnt += n + 1; // the top bit of 2^x .. old n
  }
  return cnt;
}
\end{lstlisting}
\subsection{XOR Basis}
\begin{lstlisting}
// Keeps a basis of the inserted numbers. O(B) per call.
struct Basis {
  static const int B = 60; // numbers < 2^(B+1)
  int basis[B + 1], sz = 0;
  Basis() { memset(basis, 0, sizeof(basis)); }
  // returns false if x was already representable
  bool insert(int x) {
    for (int i = B; i >= 0; i--) {
      if (!((x >> i) & 1)) continue;
      if (!basis[i]) { basis[i] = x; sz++; return true; }
      x ^= basis[i];
    }
    return false;
  }
  // can x be written as XOR of some inserted numbers?
  bool check(int x) {
    for (int i = B; i >= 0; i--) {
      if (!((x >> i) & 1)) continue;
      if (!basis[i]) return false;
      x ^= basis[i];
    }
    return true;
  }
  int maxXor() { // maximum XOR of any subset
    int res = 0;
    for (int i = B; i >= 0; i--)
      if ((res ^ basis[i]) > res) res ^= basis[i];
    return res;
  }
  // number of different subset XOR values = 2^sz
};
\end{lstlisting}
\section{Constructive}
\subsection{Convert base}
\begin{lstlisting}
// Digits 0-9 then A-Z (base <= 36). Value fits in ll.
int digitVal(char c) {
  return isdigit(c) ? c - '0' : c - 'A' + 10;
}
char digitChar(int v) {
  return v < 10 ? char('0' + v) : char('A' + v - 10);
}
int toDeci(const string &s, int base) {
  int num = 0;
  for (char c : s) num = num * base + digitVal(c);
  return num;
}
string fromDeci(int num, int base) {
  if (num == 0) return "0";
  string res;
  while (num > 0) {
    res += digitChar(num % base);
    num /= base;
  }
  reverse(all(res));
  return res;
}
// number s written in base a -> same number in base b
string convertBase(const string &s, int a, int b) {
  return fromDeci(toDeci(s, a), b);
}
\end{lstlisting}
\subsection{Fraction}
\begin{lstlisting}
// Exact a/b arithmetic, always reduced, with den > 0.
// num * den of two fractions must fit in ll.
struct Fraction {
  int num, den;
  Fraction(int n = 0, int d = 1) : num(n), den(d) {
    if (den < 0) num = -num, den = -den;
    int g = __gcd(abs(num), den);
    if (g) num /= g, den /= g;
  }
  Fraction operator+(const Fraction &o) const {
    return Fraction(num * o.den + o.num * den,
                    den * o.den);
  }
  Fraction operator-(const Fraction &o) const {
    return Fraction(num * o.den - o.num * den,
                    den * o.den);
  }
  Fraction operator*(const Fraction &o) const {
    return Fraction(num * o.num, den * o.den);
  }
  Fraction operator/(const Fraction &o) const {
    return Fraction(num * o.den, den * o.num); // o != 0
  }
  bool operator<(const Fraction &o) const {
    return num * o.den < o.num * den;
  }
  bool operator==(const Fraction &o) const {
    return num == o.num && den == o.den;
  }
  double value() const { return (double)num / den; }
};
// print: cout << f.num << "/" << f.den;
\end{lstlisting}
\section{Number Theory}
\subsection{Mod power and inverse}
\begin{lstlisting}
// (a^b) % m in O(log b). Needs m*m to fit in ll (m<=3e9)
int power(int a, int b, int m = MOD) {
  int res = 1 % m;
  a %= m;
  if (a < 0) a += m;
  while (b > 0) {
    if (b & 1) res = res * a % m;
    a = a * a % m;
    b >>= 1;
  }
  return res;
}
// inverse of a when m is PRIME (Fermat): a^(m-2)
int modInv(int a, int m = MOD) {
  return power(a, m - 2, m);
}
// a^e % m where e is a huge decimal number in a string
int powStr(int a, const string &e, int m = MOD) {
  int res = 1 % m;
  for (char c : e) // res = res^10 * a^digit
    res = power(res, 10, m) * power(a, c - '0', m) % m;
  return res;
}
\end{lstlisting}
\subsection{Extended Euclid and mod inverse}
\begin{lstlisting}
// returns g = gcd(a, b) and x, y with a*x + b*y = g
int extgcd(int a, int b, int &x, int &y) {
  if (b == 0) { x = 1; y = 0; return a; }
  int x1, y1;
  int g = extgcd(b, a % b, x1, y1);
  x = y1;
  y = x1 - y1 * (a / b);
  return g;
}
// inverse of a mod m, m does NOT need to be prime.
// returns -1 when gcd(a, m) != 1 (no inverse)
int modInverse(int a, int m) {
  int x, y;
  a = (a % m + m) % m;
  if (extgcd(a, m, x, y) != 1) return -1;
  return (x % m + m) % m;
}
// a*x + b*y = c has a solution <=> c % gcd(a, b) == 0.
// one solution: x * (c/g), y * (c/g); all solutions:
// x + k * (b/g), y - k * (a/g)
\end{lstlisting}
\subsection{Chinese Remainder Theorem}
\begin{lstlisting}
// x = a1 (mod m1) and x = a2 (mod m2); m1, m2 need NOT
// be coprime. Returns {x, lcm(m1,m2)} or {-1,-1} if no
// solution. lcm must fit in ll. Needs extgcd().
pair<int, int> crt(int a1, int m1, int a2, int m2) {
  int p, q;
  int g = extgcd(m1, m2, p, q);
  if ((a2 - a1) % g != 0) return {-1, -1};
  int l = m1 / g * m2, md = m2 / g;
  int k = (__int128)((a2 - a1) / g % md) * p % md;
  int x = (a1 + (__int128)m1 * k) % l;
  if (x < 0) x += l;
  return {x, l};
}
// Many equations: fold them one by one.
//  pair<int,int> r = {a[0], m[0]};
//  for (i = 1..k-1)
//    r = crt(r.first, r.second, a[i], m[i]);
\end{lstlisting}
\subsection{Mod inverse 1 to N}
\begin{lstlisting}
// inv[i] for all 1 <= i <= n in O(n). MOD prime, n < MOD
const int N = 1e6 + 5;
int inv[N];
void preInv(int n) {
  inv[1] = 1;
  for (int i = 2; i <= n; i++)
    inv[i] = (MOD - MOD / i * inv[MOD % i] % MOD) % MOD;
}
\end{lstlisting}
\subsection{nCr, nPr}
\begin{lstlisting}
// Needs power(). Call preFact() once before any query.
const int N = 1e6 + 5; // CHANGE: max n + 1
int fact[N], ifact[N];
void preFact() { // O(N)
  fact[0] = 1;
  for (int i = 1; i < N; i++)
    fact[i] = fact[i - 1] * i % MOD;
  ifact[N - 1] = power(fact[N - 1], MOD - 2);
  for (int i = N - 1; i >= 1; i--)
    ifact[i - 1] = ifact[i] * i % MOD;
}
int nCr(int n, int r) { // O(1)
  if (r < 0 || r > n) return 0;
  return fact[n] * ifact[r] % MOD * ifact[n - r] % MOD;
}
int nPr(int n, int r) {
  if (r < 0 || r > n) return 0;
  return fact[n] * ifact[n - r] % MOD;
}
// Small n, or MOD not prime: Pascal triangle, O(n^2)
int C[405][405];
void prePascal() {
  for (int i = 0; i <= 400; i++) {
    C[i][0] = 1;
    for (int j = 1; j <= i; j++)
      C[i][j] = (C[i - 1][j] + C[i - 1][j - 1]) % MOD;
  }
}
// Catalan(n) = nCr(2n, n) / (n + 1)
// Derangement: D(n) = (n - 1) * (D(n-1) + D(n-2))
\end{lstlisting}
\subsection{Lucas and Legendre}
\begin{lstlisting}
// Lucas: nCr % p for HUGE n, r and a SMALL prime p.
// Setup: set MOD = p and N = p (tables hold 0..p-1),
// call preFact(), then lucas(n, r). O(p + log_p n)
int lucas(int n, int r) {
  if (r < 0 || r > n) return 0;
  int res = 1;
  while (n > 0 || r > 0) { // digit by digit in base p
    res = res * nCr(n % MOD, r % MOD) % MOD;
    n /= MOD;
    r /= MOD;
  }
  return res;
}
// Legendre: largest e such that p^e divides n! (p prime)
int legendre(int n, int p) {
  int e = 0;
  while (n) {
    n /= p;
    e += n;
  }
  return e;
}
// trailing zeros of n! = legendre(n, 5)
\end{lstlisting}
\subsection{Sieve (linear, smallest prime factor)}
\begin{lstlisting}
// Memory: int is 8 bytes here. For N = 1e7 declare the
// array as int32_t spf[N] to use 40 MB instead of 80 MB.
const int N = 1e6 + 5;
int spf[N]; // spf[x] = smallest prime factor of x
vector<int> primes;
void sieve() { // O(N). x >= 2 is prime <=> spf[x] == x
  for (int i = 2; i < N; i++) {
    if (spf[i] == 0) {
      spf[i] = i;
      primes.push_back(i);
    }
    for (int p : primes) {
      if (p > spf[i] || i * p >= N) break;
      spf[i * p] = p;
    }
  }
}
// {prime, power} list of x < N in O(log x)
vector<pair<int, int>> factorize(int x) {
  vector<pair<int, int>> f;
  while (x > 1) {
    int p = spf[x], c = 0;
    while (x % p == 0) x /= p, c++;
    f.push_back({p, c});
  }
  return f;
}
\end{lstlisting}
\subsection{Bitset prime sieve (up to 1e8)}
\begin{lstlisting}
// ~12 MB for the bitset. Only odd numbers are crossed.
const int MAX = 1e8;
bitset<MAX + 1> is_prime;
vector<int32_t> primes; // 5761455 primes below 1e8
void fastSieve() {
  is_prime.set();
  is_prime[0] = is_prime[1] = 0;
  for (int i = 4; i <= MAX; i += 2) is_prime[i] = 0;
  for (int i = 3; i * i <= MAX; i += 2)
    if (is_prime[i])
      for (int j = i * i; j <= MAX; j += i * 2)
        is_prime[j] = 0;
  primes.push_back(2);
  for (int i = 3; i <= MAX; i += 2)
    if (is_prime[i]) primes.push_back(i);
}
\end{lstlisting}
\subsection{Prime factorization and divisors}
\begin{lstlisting}
// {prime, power} list of n in O(sqrt(n))
vector<pair<int, int>> primeFactors(int n) {
  vector<pair<int, int>> f;
  for (int d = 2; d * d <= n; d++) {
    if (n % d) continue;
    int c = 0;
    while (n % d == 0) n /= d, c++;
    f.push_back({d, c});
  }
  if (n > 1) f.push_back({n, 1});
  return f;
}
// The 3 functions below take ANY {prime, power} list:
// primeFactors(n), factorize(n) or factorizeBig(n).
// n = p1^e1 * p2^e2 ... -> (e1 + 1) * (e2 + 1) ...
int numberOfDivisors(const vector<pair<int, int>> &f) {
  int total = 1;
  for (auto [p, e] : f) total *= e + 1;
  return total;
}
// product of (1 + p + p^2 + ... + p^e)
int sumOfDivisors(const vector<pair<int, int>> &f) {
  int total = 1;
  for (auto [p, e] : f) {
    int sum = 1, pw = 1;
    for (int i = 0; i < e; i++) pw *= p, sum += pw;
    total *= sum;
  }
  return total;
}
// every divisor, sorted. O(D log D)
vector<int> divisors(const vector<pair<int, int>> &f) {
  vector<int> d = {1};
  for (auto [p, e] : f) {
    int sz = d.size(), pw = 1;
    for (int i = 0; i < e; i++) {
      pw *= p;
      for (int j = 0; j < sz; j++)
        d.push_back(d[j] * pw);
    }
  }
  sort(all(d));
  return d;
}
// Divisor count of EVERY x <= n in O(n log n):
//  for (i = 1..n)
//    for (j = i; j <= n; j += i) cntDiv[j]++;
\end{lstlisting}
\subsection{Euler Totient}
\begin{lstlisting}
// phi(n) = how many 1 <= x <= n have gcd(x, n) == 1
const int N = 1e6 + 5;
int phi[N];
void phiSieve() { // all phi[0..N-1] in O(N log log N)
  for (int i = 0; i < N; i++) phi[i] = i;
  for (int i = 2; i < N; i++)
    if (phi[i] == i) // i is prime
      for (int j = i; j < N; j += i)
        phi[j] -= phi[j] / i;
}
int phiSingle(int n) { // one value in O(sqrt(n))
  int res = n;
  for (int i = 2; i * i <= n; i++) {
    if (n % i) continue;
    while (n % i == 0) n /= i;
    res -= res / i;
  }
  if (n > 1) res -= res / n;
  return res;
}
// phi of every number in [L, R]. R <= 1e12, R - L <= 1e6
vector<int> phiRange(int L, int R) {
  int len = R - L + 1, lim = sqrtl(R) + 1;
  vector<int> ph(len), rem(len);
  for (int i = 0; i < len; i++) ph[i] = rem[i] = L + i;
  vector<bool> comp(lim + 1, false);
  for (int p = 2; p <= lim; p++) {
    if (comp[p]) continue;
    for (int j = p * p; j <= lim; j += p) comp[j] = true;
    for (int j = max(p, (L + p - 1) / p) * p; j <= R;
         j += p) {
      ph[j - L] -= ph[j - L] / p;
      while (rem[j - L] % p == 0) rem[j - L] /= p;
    }
  }
  for (int i = 0; i < len; i++) // one big prime left
    if (rem[i] > 1) ph[i] -= ph[i] / rem[i];
  return ph;
}
// Euler: a^phi(m) = 1 (mod m) when gcd(a, m) == 1
//  so inverse of a mod m = power(a, phiSingle(m) - 1, m)
// Huge exponent b >= phi(m), any a:
//  a^b = a^(b % phi(m) + phi(m)) (mod m)
// sum of phi(d) over all divisors d of n = n
\end{lstlisting}
\subsection{Mobius and coprime pairs}
\begin{lstlisting}
// mu[1] = 1; mu[n] = 0 if a square divides n;
// otherwise (-1)^(number of prime factors). O(n log n)
vector<int> mobius(int n) {
  vector<int> mu(n + 1, 1), isp(n + 1, 1);
  mu[0] = 0;
  for (int i = 2; i <= n; i++) {
    if (!isp[i]) continue;
    for (int j = i; j <= n; j += i) {
      mu[j] *= -1;
      if (j > i) isp[j] = 0;
    }
    for (int j = i * i; j <= n; j += i * i) mu[j] = 0;
  }
  return mu;
}
// pairs i < j with gcd(a[i], a[j]) == 1. O(A log A)
int coprimePairs(const vector<int> &a) {
  int mx = *max_element(all(a));
  vector<int> freq(mx + 1, 0), mu = mobius(mx);
  for (int x : a) freq[x]++;
  int res = 0;
  for (int k = 1; k <= mx; k++) {
    if (mu[k] == 0) continue;
    int d = 0; // how many elements are multiples of k
    for (int j = k; j <= mx; j += k) d += freq[j];
    res += mu[k] * (d * (d - 1) / 2);
  }
  return res;
}
\end{lstlisting}
\subsection{Miller Rabin and Pollard Rho}
\begin{lstlisting}
// isPrime: exact for every n < 2^63, O(12 * log n).
// factorizeBig: factorization of n <= 1e18, ~O(n^(1/4)).
// Uses rng from the template.
ll mulmod(ll a, ll b, ll m) {
  return (__int128)a * b % m;
}
ll powmod(ll a, ll b, ll m) {
  ll res = 1;
  a %= m;
  while (b > 0) {
    if (b & 1) res = mulmod(res, a, m);
    a = mulmod(a, a, m);
    b >>= 1;
  }
  return res;
}
bool isPrime(ll n) {
  if (n < 2) return false;
  static const vector<ll> bases = {
    2, 3, 5, 7, 11, 13, 17, 19, 23, 29, 31, 37};
  for (ll p : bases)
    if (n % p == 0) return n == p;
  ll d = n - 1;
  int s = 0;
  while (d % 2 == 0) d /= 2, s++; // n - 1 = d * 2^s
  for (ll a : bases) {
    ll x = powmod(a, d, n);
    if (x == 1 || x == n - 1) continue;
    bool composite = true;
    for (int r = 1; r < s && composite; r++) {
      x = mulmod(x, x, n);
      if (x == n - 1) composite = false;
    }
    if (composite) return false;
  }
  return true;
}
// a non-trivial factor of a COMPOSITE n
ll pollardRho(ll n) {
  if (n % 2 == 0) return 2;
  while (true) {
    ll c = rng() % (n - 1) + 1, x = rng() % n, y = x;
    ll d = 1;
    auto f = [&](ll v) {
      return (ll)(((__int128)v * v + c) % n);
    };
    while (d == 1) { // Floyd: y moves twice as fast
      x = f(x);
      y = f(f(y));
      d = __gcd(x > y ? x - y : y - x, n);
    }
    if (d != n) return d; // d == n: retry with new c
  }
}
void factorRec(ll n, vector<ll> &f) {
  if (n == 1) return;
  if (isPrime(n)) { f.push_back(n); return; }
  ll d = pollardRho(n);
  factorRec(d, f);
  factorRec(n / d, f);
}
// sorted {prime, power} list
vector<pair<int, int>> factorizeBig(ll n) {
  vector<ll> f;
  factorRec(n, f);
  sort(all(f));
  vector<pair<int, int>> res;
  for (ll p : f) {
    if (res.empty() || res.back().first != p)
      res.push_back({p, 0});
    res.back().second++;
  }
  return res;
}
\end{lstlisting}
\subsection{Baby Step Giant Step (discrete log)}
\begin{lstlisting}
// Smallest x >= 0 with a^x = b (mod m), or -1 if none.
// m can be ANY number (not only prime). O(sqrt(m)).
// Needs power(). For x >= 1 only: handle a^0 yourself.
int bsgs(int a, int b, int m) {
  a %= m, b %= m;
  if (m == 1 || b == 1) return 0;
  // remove common factors so that gcd(a, m) == 1
  int k = 1, add = 0, g;
  while ((g = __gcd(a, m)) > 1) {
    if (b == k) return add;
    if (b % g) return -1;
    b /= g, m /= g, add++;
    k = k * (a / g) % m;
  }
  // now solve k * a^(n*p) = b * a^q, answer n*p - q
  int n = sqrtl(m) + 1;
  unordered_map<int, int> vals; // b * a^q -> largest q
  for (int q = 0, cur = b; q <= n; q++) {
    vals[cur] = q;
    cur = cur * a % m;
  }
  int an = power(a, n, m);
  for (int p = 1, cur = k; p <= n; p++) {
    cur = cur * an % m;
    if (vals.count(cur)) return n * p - vals[cur] + add;
  }
  return -1;
}
\end{lstlisting}
\subsection{Matrix multiplication and exponentiation}
\begin{lstlisting}
typedef vector<vector<int>> Mat;
// (n x k) * (k x m) matrix product mod MOD. O(n*k*m)
Mat matMul(const Mat &a, const Mat &b) {
  int n = a.size(), k = b.size(), m = b[0].size();
  Mat res(n, vector<int>(m, 0));
  for (int i = 0; i < n; i++)
    for (int x = 0; x < k; x++) {
      if (a[i][x] == 0) continue;
      for (int j = 0; j < m; j++)
        res[i][j] =
          (res[i][j] + a[i][x] * b[x][j]) % MOD;
    }
  return res;
}
// a^p for a square matrix. O(n^3 log p)
Mat matPow(Mat a, int p) {
  int n = a.size();
  Mat res(n, vector<int>(n, 0));
  for (int i = 0; i < n; i++) res[i][i] = 1; // identity
  while (p > 0) {
    if (p & 1) res = matMul(res, a);
    a = matMul(a, a);
    p >>= 1;
  }
  return res;
}
// Usage for f(n) = c1*f(n-1) + c2*f(n-2):
//  T = {{c1, c2}, {1, 0}};
//  [f(n), f(n-1)] = matPow(T, n - 1) * [f(1), f(0)]
// Fibonacci: matPow({{1,1},{1,0}}, n)[0][1] == F(n)
\end{lstlisting}
\subsection{Gaussian Elimination}
\begin{lstlisting}
// Solves A x = B. a = n rows of [A | B] (m + 1 columns).
// Returns 0 (no solution), 1 (unique) or 2 (infinite);
// ans = one solution. O(n * m * min(n, m))
const double EPS = 1e-9;
int gauss(vector<vector<double>> a,
          vector<double> &ans) {
  int n = a.size(), m = a[0].size() - 1;
  vector<int> where(m, -1); // row that fixes variable i
  for (int col = 0, row = 0; col < m && row < n; ++col) {
    int sel = row; // biggest pivot for stability
    for (int i = row; i < n; ++i)
      if (abs(a[i][col]) > abs(a[sel][col])) sel = i;
    if (abs(a[sel][col]) < EPS) continue; // free var
    swap(a[sel], a[row]);
    where[col] = row;
    for (int i = 0; i < n; ++i) {
      if (i == row) continue;
      double c = a[i][col] / a[row][col];
      for (int j = col; j <= m; ++j)
        a[i][j] -= a[row][j] * c;
    }
    ++row;
  }
  ans.assign(m, 0);
  for (int i = 0; i < m; ++i)
    if (where[i] != -1)
      ans[i] = a[where[i]][m] / a[where[i]][i];
  for (int i = 0; i < n; ++i) { // check every equation
    double sum = 0;
    for (int j = 0; j < m; ++j) sum += ans[j] * a[i][j];
    if (abs(sum - a[i][m]) > EPS) return 0;
  }
  for (int i = 0; i < m; ++i)
    if (where[i] == -1) return 2;
  return 1;
}
\end{lstlisting}
\subsection{Determinant modulo prime}
\begin{lstlisting}
// det of an n x n matrix modulo the prime MOD. O(n^3).
// Entries must be in [0, MOD). Needs power().
int determinant(vector<vector<int>> a) {
  int n = a.size(), det = 1;
  for (int col = 0; col < n; col++) {
    int piv = col; // row with a non-zero in this column
    while (piv < n && a[piv][col] == 0) piv++;
    if (piv == n) return 0;
    if (piv != col) {
      swap(a[piv], a[col]);
      det = (MOD - det) % MOD; // swap flips the sign
    }
    det = det * a[col][col] % MOD;
    int inv = power(a[col][col], MOD - 2);
    for (int r = col + 1; r < n; r++) {
      int f = a[r][col] * inv % MOD;
      for (int c = col; c < n; c++)
        a[r][c] =
          ((a[r][c] - f * a[col][c]) % MOD + MOD) % MOD;
    }
  }
  return det;
}
// Solving A x = b modulo a prime: the same steps as in
// Gaussian Elimination, but "divide by x" becomes
// "multiply by power(x, MOD - 2)".
// XOR equations (mod 2): keep rows as bitset and use
// row[i] ^= row[pivot].
\end{lstlisting}
\subsection{FFT}
\begin{lstlisting}
typedef complex<double> cd;
const double PI = acos(-1.0);
// a.size() must be a power of two
void fft(vector<cd> &a, bool invert) {
  int n = a.size();
  for (int i = 1, j = 0; i < n; i++) { // bit reversal
    int bit = n >> 1;
    for (; j & bit; bit >>= 1) j ^= bit;
    j ^= bit;
    if (i < j) swap(a[i], a[j]);
  }
  for (int len = 2; len <= n; len <<= 1) {
    double ang = 2 * PI / len * (invert ? -1 : 1);
    cd wlen(cos(ang), sin(ang));
    for (int i = 0; i < n; i += len) {
      cd w(1);
      for (int j = 0; j < len / 2; j++) {
        cd u = a[i + j], v = a[i + j + len / 2] * w;
        a[i + j] = u + v;
        a[i + j + len / 2] = u - v;
        w *= wlen;
      }
    }
  }
  if (invert)
    for (cd &x : a) x /= n;
}
// c[k] = sum of a[i] * b[k - i]. O(n log n).
// Exact while every c[k] < ~1e14 (double precision).
vector<ll> multiply(const vector<int> &a,
                    const vector<int> &b) {
  vector<cd> fa(all(a)), fb(all(b));
  int n = 1, need = a.size() + b.size() - 1;
  while (n < need) n <<= 1;
  fa.resize(n), fb.resize(n);
  fft(fa, false), fft(fb, false);
  for (int i = 0; i < n; i++) fa[i] *= fb[i];
  fft(fa, true);
  vector<ll> res(need);
  for (int i = 0; i < need; i++)
    res[i] = llround(fa[i].real());
  return res;
}
\end{lstlisting}
\subsection{NTT}
\begin{lstlisting}
// Exact multiplication modulo 998244353. Needs power().
// Other NTT primes (root 3): 469762049, 167772161
const int NMOD = 998244353, ROOT = 3;
void ntt(vector<int> &a, bool invert) {
  int n = a.size();
  for (int i = 1, j = 0; i < n; i++) {
    int bit = n >> 1;
    for (; j & bit; bit >>= 1) j ^= bit;
    j ^= bit;
    if (i < j) swap(a[i], a[j]);
  }
  for (int len = 2; len <= n; len <<= 1) {
    int wlen = power(ROOT, (NMOD - 1) / len, NMOD);
    if (invert) wlen = power(wlen, NMOD - 2, NMOD);
    for (int i = 0; i < n; i += len) {
      int w = 1;
      for (int j = 0; j < len / 2; j++) {
        int u = a[i + j];
        int v = a[i + j + len / 2] * w % NMOD;
        a[i + j] = (u + v) % NMOD;
        a[i + j + len / 2] = (u - v + NMOD) % NMOD;
        w = w * wlen % NMOD;
      }
    }
  }
  if (invert) {
    int ninv = power(n, NMOD - 2, NMOD);
    for (int &x : a) x = x * ninv % NMOD;
  }
}
// values of a, b must already be in [0, NMOD)
vector<int> multiplyMod(vector<int> a, vector<int> b) {
  int n = 1, need = a.size() + b.size() - 1;
  while (n < need) n <<= 1;
  a.resize(n), b.resize(n);
  ntt(a, false), ntt(b, false);
  for (int i = 0; i < n; i++) a[i] = a[i] * b[i] % NMOD;
  ntt(a, true);
  a.resize(need);
  return a;
}
// polynomial a^k: binary exponentiation with multiplyMod
\end{lstlisting}
\subsection{Big integer (string) multiply and power}
\begin{lstlisting}
// a * b for decimal strings. O(|a| * |b|)
string mulStr(const string &a, const string &b) {
  int n = a.size(), m = b.size();
  vector<int> r(n + m, 0);
  for (int i = n - 1; i >= 0; --i) {
    int carry = 0;
    for (int j = m - 1; j >= 0; --j) {
      int sum = r[i + j + 1] +
                (a[i] - '0') * (b[j] - '0') + carry;
      r[i + j + 1] = sum % 10;
      carry = sum / 10;
    }
    r[i] += carry;
  }
  string s;
  for (int x : r)
    if (!(s.empty() && x == 0)) s += char('0' + x);
  return s.empty() ? "0" : s;
}
// exact x^n as a decimal string
string bigPow(int x, int n) {
  string res = "1", base = to_string(x);
  while (n > 0) {
    if (n & 1) res = mulStr(res, base);
    base = mulStr(base, base);
    n >>= 1;
  }
  return res;
}
\end{lstlisting}
\subsection{Divisibility rules}
\begin{lstlisting}[language={}]
 2: last digit is even
 3: sum of digits divisible by 3
 4: last two digits divisible by 4
 5: last digit is 0 or 5
 6: divisible by both 2 and 3
 7: double the last digit, subtract it from the rest;
    repeat. 3976 -> 397 - 12 = 385 -> 38 - 10 = 28 = 7*4
 8: last three digits divisible by 8
 9: sum of digits divisible by 9
10: last digit is 0
11: (sum of digits at odd positions) - (sum at even
    positions) divisible by 11.
    8153783 -> (8+5+7+3) - (1+3+8) = 11  -> divisible
12: divisible by both 3 and 4
13: 4 * last digit + the rest; repeat.
    585 -> 58 + 20 = 78 = 13 * 6
    Long numbers: split into groups of 3 digits from the
    right; alternating sum of groups divisible by 13
    (the same trick works for 7 and 11).
    1108503773 -> (773 + 108) - (503 + 1) = 377 = 13*29
15: divisible by both 3 and 5
\end{lstlisting}
\subsection{Formulas}
\begin{lstlisting}[language={}]
Sums
 1 + 2 + ... + n       = n(n+1)/2
 1^2 + 2^2 + ... + n^2 = n(n+1)(2n+1)/6
 1^3 + 2^3 + ... + n^3 = (n(n+1)/2)^2
 a + ar + ... + ar^(n-1) = a(r^n - 1)/(r - 1)
Number theory
 gcd(a, b) * lcm(a, b) = a * b
 Fermat: a^(p-1) = 1 (mod p), p prime, p not dividing a
 Wilson: (p-1)! = -1 (mod p)  <=>  p is prime
 (a / b) mod p = a * b^(p-2) mod p
 primes below 1e6: 78498, 1e7: 664579, 1e8: 5761455
 max number of divisors: n <= 1e6: 240, 1e9: 1344,
  1e12: 6720, 1e18: 103680
 gcd(F(m), F(n)) = F(gcd(m, n));
 F(1) + ... + F(n) = F(n+2) - 1
Combinatorics
 C(n,r) = C(n-1,r) + C(n-1,r-1) = C(n,n-r)
 sum of C(n,i) = 2^n;  sum of i*C(n,i) = n * 2^(n-1)
 hockey stick: C(r,r) + C(r+1,r) + ... + C(n,r)
               = C(n+1,r+1)
 Vandermonde: sum of C(m,i)*C(n,k-i) = C(m+n,k)
 Catalan 1 1 2 5 14 42 132 429 1430: balanced brackets,
  binary trees, triangulations, paths below diagonal
 Derangements 1 0 1 2 9 44 265 (nobody gets own item)
 n items in a circle: (n-1)! orders
 inclusion-exclusion: |A or B| = |A| + |B| - |A and B|
Geometry
 triangle area = sqrt(s(s-a)(s-b)(s-c)), s = (a+b+c)/2
 circumradius R = abc / (4 * area); inradius r = area / s
 c^2 = a^2 + b^2 - 2ab*cos(C);  a / sin(A) = 2R
 Pick: area = inside + border/2 - 1
 Euler (planar graph): V - E + F = 2
 sphere: volume 4/3*PI*r^3, surface 4*PI*r^2
Probability
 E[X + Y] = E[X] + E[Y] always (linearity)
 expected tries until success with chance p = 1/p
\end{lstlisting}
\section{String}
\subsection{Hashing}
\begin{lstlisting}
// Double polynomial hash. Indices are 0-based, [l, r]
// inclusive. Build O(n), every query O(1).
// Usage: PolyHash h(s); h.get(l, r) == h.get(x, y)
struct PolyHash {
  static const ll M1 = 1e9 + 7, M2 = 1e9 + 9;
  static const ll P1 = 257, P2 = 263;
  static vector<ll> pw1, pw2; // powers of the base
  vector<ll> h1, h2, r1, r2; // hash of s / reversed s
  int n;
  static void grow(int len) {
    while ((int)pw1.size() <= len) {
      pw1.push_back(pw1.back() * P1 % M1);
      pw2.push_back(pw2.back() * P2 % M2);
    }
  }
  PolyHash(const string &s) : n(s.size()) {
    grow(n);
    h1.assign(n + 1, 0), h2 = r1 = r2 = h1;
    for (int i = 0; i < n; ++i) {
      h1[i + 1] = (h1[i] * P1 + s[i]) % M1;
      h2[i + 1] = (h2[i] * P2 + s[i]) % M2;
      r1[i + 1] = (r1[i] * P1 + s[n - 1 - i]) % M1;
      r2[i + 1] = (r2[i] * P2 + s[n - 1 - i]) % M2;
    }
  }
  // hash of s[l..r]
  pair<ll, ll> get(int l, int r) {
    if (l > r) return {0, 0};
    int len = r - l + 1;
    ll a = (h1[r + 1] - h1[l] * pw1[len] % M1 + M1) % M1;
    ll b = (h2[r + 1] - h2[l] * pw2[len] % M2 + M2) % M2;
    return {a, b};
  }
  // hash of s[l..r] read backwards (from r down to l)
  pair<ll, ll> getRev(int l, int r) {
    if (l > r) return {0, 0};
    int len = r - l + 1, L = n - 1 - r, R = n - 1 - l;
    ll a = (r1[R + 1] - r1[L] * pw1[len] % M1 + M1) % M1;
    ll b = (r2[R + 1] - r2[L] * pw2[len] % M2 + M2) % M2;
    return {a, b};
  }
  bool isPalindrome(int l, int r) {
    return get(l, r) == getRev(l, r);
  }
  // hash of (string A + string B) from their hashes
  static pair<ll, ll> concat(pair<ll, ll> hA,
                             pair<ll, ll> hB, int lenB) {
    grow(lenB);
    return {(hA.first * pw1[lenB] + hB.first) % M1,
            (hA.second * pw2[lenB] + hB.second) % M2};
  }
  // longest common prefix of suffixes i and j. O(log n)
  int lcp(int i, int j) {
    int lo = 0, hi = n - max(i, j);
    while (lo < hi) {
      int mid = (lo + hi + 1) / 2;
      if (get(i, i + mid - 1) == get(j, j + mid - 1))
        lo = mid;
      else hi = mid - 1;
    }
    return lo;
  }
  // compare s[a..b] with s[c..d]: -1, 0 or 1
  // (pass the string s itself). O(log n)
  int compare(const string &s, int a, int b, int c,
              int d) {
    int len1 = b - a + 1, len2 = d - c + 1;
    int k = min({lcp(a, c), len1, len2});
    if (k == min(len1, len2))
      return len1 < len2 ? -1 : (len1 > len2 ? 1 : 0);
    return s[a + k] < s[c + k] ? -1 : 1;
  }
  // does s[l..r] repeat with period p ?
  bool isPeriodic(int l, int r, int p) {
    return get(l, r - p) == get(l + p, r);
  }
};
vector<ll> PolyHash::pw1{1};
vector<ll> PolyHash::pw2{1};
\end{lstlisting}
\subsection{KMP Algorithm}
\begin{lstlisting}
// pi[i] = length of the longest proper prefix of
// p[0..i] that is also a suffix of it. O(n + m)
struct KMP {
  string pat;
  vector<int> pi;
  KMP(const string &p) : pat(p) {
    int m = pat.size();
    pi.assign(m, 0);
    for (int i = 1, j = 0; i < m; i++) {
      while (j > 0 && pat[i] != pat[j]) j = pi[j - 1];
      if (pat[i] == pat[j]) j++;
      pi[i] = j;
    }
  }
  // start index of every occurrence of pat in text
  vector<int> search(const string &text) {
    vector<int> res;
    int n = text.size(), m = pat.size();
    if (m == 0) return res;
    for (int i = 0, j = 0; i < n; i++) {
      while (j > 0 && text[i] != pat[j]) j = pi[j - 1];
      if (text[i] == pat[j]) j++;
      if (j == m) {
        res.push_back(i - m + 1);
        j = pi[j - 1];
      }
    }
    return res;
  }
};
// smallest period of s = n - pi[n - 1]
\end{lstlisting}
\subsection{Z Algorithm}
\begin{lstlisting}
// z[i] = longest common prefix of s and s[i..]. O(n)
vector<int> zFunction(const string &s) {
  int n = s.size();
  vector<int> z(n, 0);
  for (int i = 1, l = 0, r = 0; i < n; ++i) {
    if (i <= r) z[i] = min(r - i + 1, z[i - l]);
    while (i + z[i] < n && s[z[i]] == s[i + z[i]])
      ++z[i];
    if (i + z[i] - 1 > r) l = i, r = i + z[i] - 1;
  }
  return z;
}
// start index of every occurrence of pat in text.
// '$' must not appear in either string.
vector<int> zSearch(const string &text,
                    const string &pat) {
  vector<int> z = zFunction(pat + "$" + text), res;
  int m = pat.size();
  for (int i = m + 1; i < (int)z.size(); i++)
    if (z[i] == m) res.push_back(i - (m + 1));
  return res;
}
\end{lstlisting}
\subsection{Manacher's Algorithm}
\begin{lstlisting}
// All palindromic substrings in O(n).
// Internally t = ^#a#b#c#$ ; p[i] = palindrome radius.
struct Manacher {
  vector<int> p;
  Manacher(const string &s) {
    string t = "^";
    for (char c : s) t += "#" + string(1, c);
    t += "#$";
    int n = t.size(), c = 0, r = 0;
    p.assign(n, 0);
    for (int i = 1; i < n - 1; i++) {
      if (i < r) p[i] = min(r - i, p[2 * c - i]);
      while (t[i + 1 + p[i]] == t[i - 1 - p[i]]) p[i]++;
      if (i + p[i] > r) c = i, r = i + p[i];
    }
  }
  // is s[l..r] a palindrome? (0-based, inclusive)
  bool isPalindrome(int l, int r) {
    return p[l + r + 2] >= r - l + 1;
  }
  // longest palindrome with center at s[i] (odd length)
  int oddLen(int i) { return p[2 * i + 2]; }
  // longest palindrome centered between s[i], s[i+1]
  int evenLen(int i) { return p[2 * i + 3]; }
  string longest(const string &s) {
    int best = 0, center = 0;
    for (int i = 1; i < (int)p.size() - 1; i++)
      if (p[i] > best) best = p[i], center = i;
    return s.substr((center - best) / 2, best);
  }
};
\end{lstlisting}
\subsection{Trie}
\begin{lstlisting}
// Lowercase words. CHANGE 26 and 'a' for other letters.
struct Trie {
  vector<array<int, 26>> nxt; // 0 = no child
  vector<int> cnt, pre; // words ending / passing a node
  Trie() { newNode(); }
  int newNode() {
    nxt.push_back({});
    cnt.push_back(0), pre.push_back(0);
    return nxt.size() - 1;
  }
  void insert(const string &s) {
    int u = 0;
    for (char c : s) {
      if (!nxt[u][c - 'a']) {
        int v = newNode();
        nxt[u][c - 'a'] = v;
      }
      u = nxt[u][c - 'a'];
      pre[u]++;
    }
    cnt[u]++;
  }
  int find(const string &s) { // node of s, or -1
    int u = 0;
    for (char c : s) {
      if (!nxt[u][c - 'a']) return -1;
      u = nxt[u][c - 'a'];
    }
    return u;
  }
  // how many inserted words are equal to s
  int count(const string &s) {
    int u = find(s);
    return u == -1 ? 0 : cnt[u];
  }
  // how many inserted words start with s
  int countPrefix(const string &s) {
    int u = find(s);
    return u == -1 ? 0 : pre[u];
  }
};
\end{lstlisting}
\subsection{Binary Trie}
\begin{lstlisting}
// Multiset of non-negative numbers for XOR queries.
// O(B) per operation. Memory: ~3 * B ints per number.
struct BinaryTrie {
  static const int B = 30; // CHANGE: numbers < 2^(B+1)
  vector<array<int, 2>> ch; // 0 = no child
  vector<int> cnt; // how many numbers pass each node
  BinaryTrie() : ch(1), cnt(1, 0) {} // node 0 = root
  // d = +1: insert x.   d = -1: erase x (x must exist)
  void add(int x, int d) {
    int u = 0;
    cnt[u] += d;
    for (int i = B; i >= 0; i--) {
      int b = (x >> i) & 1;
      if (!ch[u][b]) {
        int v = ch.size();
        ch.push_back({0, 0}), cnt.push_back(0);
        ch[u][b] = v;
      }
      u = ch[u][b];
      cnt[u] += d;
    }
  }
  int size() { return cnt[0]; }
  // max of (x ^ y) over stored y; -1 if empty
  int maxXor(int x) {
    if (cnt[0] == 0) return -1;
    int u = 0, ans = 0;
    for (int i = B; i >= 0; i--) {
      int b = (x >> i) & 1, want = ch[u][b ^ 1];
      if (want && cnt[want] > 0) {
        ans |= (1LL << i);
        u = want;
      } else u = ch[u][b];
    }
    return ans;
  }
  // k-th smallest (k = 1, 2, ...) value of (x ^ y);
  // -1 if k > size. min XOR = kthXor(x, 1),
  // k-th largest = kthXor(x, size() - k + 1)
  int kthXor(int x, int k) {
    if (k > cnt[0]) return -1;
    int u = 0, ans = 0;
    for (int i = B; i >= 0; i--) {
      int b = (x >> i) & 1, same = ch[u][b];
      int c = same ? cnt[same] : 0;
      if (k <= c) u = same; // this bit of the XOR is 0
      else {
        k -= c;
        ans |= (1LL << i);
        u = ch[u][b ^ 1];
      }
    }
    return ans;
  }
  // how many stored y have (x ^ y) < lim
  int countLess(int x, int lim) {
    int u = 0, res = 0;
    for (int i = B; i >= 0; i--) {
      int b = (x >> i) & 1;
      if ((lim >> i) & 1) {
        // XOR bit 0 here -> everything below is smaller
        if (ch[u][b]) res += cnt[ch[u][b]];
        u = ch[u][b ^ 1];
      } else u = ch[u][b];
      if (!u) break; // fell out of the trie
    }
    return res;
  }
};
\end{lstlisting}
\subsection{Aho Corasick}
\begin{lstlisting}
// Many patterns, one text. Build O(total length * 26).
// Usage: AhoCorasick ac; ac.insert(p) for every pattern;
// ac.build(); then countAll(text) or occurrences(text).
struct AhoCorasick {
  vector<array<int, 26>> nxt; // automaton transitions
  vector<int> fail, out, order, endNode;
  AhoCorasick() { newNode(); }
  int newNode() {
    nxt.push_back({});
    fail.push_back(0), out.push_back(0);
    return nxt.size() - 1;
  }
  void insert(const string &s) {
    int u = 0;
    for (char c : s) {
      if (!nxt[u][c - 'a']) {
        int v = newNode();
        nxt[u][c - 'a'] = v;
      }
      u = nxt[u][c - 'a'];
    }
    out[u]++;
    endNode.push_back(u); // node of the i-th pattern
  }
  void build() { // BFS: call once after all inserts
    queue<int> q;
    for (int i = 0; i < 26; i++)
      if (nxt[0][i]) q.push(nxt[0][i]);
    while (!q.empty()) {
      int u = q.front();
      q.pop();
      order.push_back(u);
      out[u] += out[fail[u]]; // patterns ending here
      for (int i = 0; i < 26; i++) {
        int v = nxt[u][i];
        if (v) {
          fail[v] = nxt[fail[u]][i];
          q.push(v);
        } else nxt[u][i] = nxt[fail[u]][i];
      }
    }
  }
  // total occurrences of all patterns in text
  int countAll(const string &text) {
    int u = 0, res = 0;
    for (char c : text) {
      u = nxt[u][c - 'a'];
      res += out[u];
    }
    return res;
  }
  // res[i] = occurrences of the i-th inserted pattern
  vector<int> occurrences(const string &text) {
    vector<int> vis(nxt.size(), 0);
    int u = 0;
    for (char c : text) vis[u = nxt[u][c - 'a']]++;
    for (int i = (int)order.size() - 1; i >= 0; i--)
      vis[fail[order[i]]] += vis[order[i]];
    vector<int> res;
    for (int e : endNode) res.push_back(vis[e]);
    return res;
  }
};
\end{lstlisting}
\subsection{Suffix Array}
\begin{lstlisting}
// Works for a string or a vector<int>. O(n log n).
// A smallest sentinel is added, so n = length + 1 and
// p[0] is the sentinel. p[1..n-1] = suffixes sorted.
// c[i] = position of suffix i inside p.
// lcp[i] = LCP of suffixes p[i-1] and p[i]  (i >= 2).
struct SuffixArray {
  int n;
  vector<int> s, p, c, lcp;
  vector<vector<int>> rmq; // sparse table over lcp
  SuffixArray(const string &t) : s(all(t)) { build(); }
  SuffixArray(const vector<int> &v) : s(v) { build(); }
  void build() {
    // compress values to 1..k, sentinel = 0
    vector<int> srt = s;
    sort(all(srt));
    srt.erase(unique(all(srt)), srt.end());
    for (int &x : s)
      x = lower_bound(all(srt), x) - srt.begin() + 1;
    s.push_back(0);
    n = s.size();
    p.resize(n), c.resize(n);
    vector<int> cnt(n + 1, 0), pn(n), cn(n);
    // sort by first character (counting sort)
    for (int i = 0; i < n; i++) cnt[s[i]]++;
    for (int i = 1; i <= n; i++) cnt[i] += cnt[i - 1];
    for (int i = n - 1; i >= 0; i--) p[--cnt[s[i]]] = i;
    int classes = 1;
    c[p[0]] = 0;
    for (int i = 1; i < n; i++) {
      if (s[p[i]] != s[p[i - 1]]) classes++;
      c[p[i]] = classes - 1;
    }
    // sort by first 2^(k+1) characters using 2^k result
    for (int len = 1; len < n; len <<= 1) {
      for (int i = 0; i < n; i++)
        pn[i] = (p[i] - len + n) % n;
      fill(cnt.begin(), cnt.begin() + classes, 0);
      for (int i = 0; i < n; i++) cnt[c[pn[i]]]++;
      for (int i = 1; i < classes; i++)
        cnt[i] += cnt[i - 1];
      for (int i = n - 1; i >= 0; i--)
        p[--cnt[c[pn[i]]]] = pn[i];
      cn[p[0]] = 0;
      classes = 1;
      for (int i = 1; i < n; i++) {
        int a = p[i], b = p[i - 1];
        pair<int, int> cur = {c[a], c[(a + len) % n]};
        pair<int, int> prv = {c[b], c[(b + len) % n]};
        if (cur != prv) classes++;
        cn[p[i]] = classes - 1;
      }
      c = cn;
    }
    // Kasai: lcp array in O(n)
    lcp.assign(n, 0);
    for (int i = 0, k = 0; i < n - 1; i++) {
      int j = p[c[i] - 1];
      while (s[i + k] == s[j + k]) k++;
      lcp[c[i]] = k;
      if (k > 0) k--;
    }
    // sparse table for range minimum on lcp
    rmq.assign(1, lcp);
    for (int j = 1; (1LL << j) <= n; j++) {
      rmq.push_back(vector<int>(n));
      for (int i = 0; i + (1LL << j) <= n; i++)
        rmq[j][i] = min(rmq[j - 1][i],
          rmq[j - 1][i + (1LL << (j - 1))]);
    }
  }
  // LCP of the suffixes starting at i and j. O(1)
  int getLcp(int i, int j) {
    if (i == j) return n - 1 - i;
    int l = c[i], r = c[j];
    if (l > r) swap(l, r);
    l++;
    int k = 63 - __builtin_clzll(r - l + 1);
    return min(rmq[k][l], rmq[k][r - (1LL << k) + 1]);
  }
  // number of different non-empty substrings
  int distinctSubstrings() {
    int total = 0;
    for (int i = 1; i < n; i++)
      total += (n - 1 - p[i]) - lcp[i];
    return total;
  }
  // occurrences of pat in t, where t is the SAME string
  // the array was built from. O(|pat| log n)
  int countPattern(const string &t, const string &pat) {
    int m = pat.size();
    auto lo = lower_bound(
      p.begin() + 1, p.end(), pat,
      [&](int i, const string &v) {
        return t.compare(i, m, v) < 0;
      });
    auto hi = upper_bound(
      p.begin() + 1, p.end(), pat,
      [&](const string &v, int i) {
        return t.compare(i, m, v) > 0;
      });
    return hi - lo;
  }
};
\end{lstlisting}
\subsection{String pattern matching (FFT)}
\begin{lstlisting}
// Counts positions where pat matches text exactly.
// Lowercase letters, O(26 * n log n). Needs multiply().
// Idea: for each letter count equal positions with FFT.
// Wildcards: skip the wildcard letter and compare with
// the number of non-wildcard characters instead of m.
int patternMatchFFT(const string &text,
                    const string &pat) {
  int n = text.size(), m = pat.size();
  if (m > n) return 0;
  vector<ll> total(n + m - 1, 0);
  for (char ch = 'a'; ch <= 'z'; ch++) {
    vector<int> A(n, 0), B(m, 0);
    for (int i = 0; i < n; i++) A[i] = (text[i] == ch);
    // reversed pattern: match ending at i lands at c[i]
    for (int i = 0; i < m; i++)
      B[m - 1 - i] = (pat[i] == ch);
    vector<ll> c = multiply(A, B);
    for (int i = 0; i < n + m - 1; i++) total[i] += c[i];
  }
  int matches = 0;
  for (int i = m - 1; i < n; i++)
    if (total[i] == m) matches++;
  return matches;
}
\end{lstlisting}
\section{DP}
\subsection{Kadane's algorithm}
\begin{lstlisting}
// Maximum subarray sum (non-empty). O(n).
// Returns {best sum, left index, right index}.
array<int, 3> kadane(const vector<int> &a) {
  int best = a[0], bl = 0, br = 0;
  int sum = 0, start = 0; // current segment = [start, r]
  for (int r = 0; r < (int)a.size(); ++r) {
    sum += a[r];
    if (sum > best) best = sum, bl = start, br = r;
    if (sum < 0) sum = 0, start = r + 1; // drop prefix
  }
  return {best, bl, br};
}
\end{lstlisting}
\subsection{LIS}
\begin{lstlisting}
// Longest STRICTLY increasing subsequence. O(n log n).
// d[l] = smallest possible last value of a length-l one.
// For non-decreasing: use upper_bound instead.
int lis(const vector<int> &a) {
  vector<int> d;
  for (int x : a) {
    auto it = lower_bound(all(d), x);
    if (it == d.end()) d.push_back(x);
    else *it = x;
  }
  return d.size();
}
// same, but also returns the subsequence itself
vector<int> lisSeq(const vector<int> &a) {
  int n = a.size();
  vector<int> d, pos, par(n, -1); // pos[l] = index in a
  for (int i = 0; i < n; i++) {
    int l = lower_bound(all(d), a[i]) - d.begin();
    if (l == (int)d.size())
      d.push_back(0), pos.push_back(0);
    d[l] = a[i], pos[l] = i;
    par[i] = l ? pos[l - 1] : -1;
  }
  vector<int> seq;
  for (int v = pos.empty() ? -1 : pos.back(); v != -1;
       v = par[v])
    seq.push_back(a[v]);
  reverse(all(seq));
  return seq;
}
\end{lstlisting}
\subsection{LCS}
\begin{lstlisting}
// Longest common subsequence. O(n * m) time and memory.
struct LCS {
  string a, b;
  int n, m;
  vector<vector<int>> dp; // dp[i][j]: a[0..i), b[0..j)
  LCS(const string &x, const string &y) : a(x), b(y) {
    n = a.size(), m = b.size();
    dp.assign(n + 1, vector<int>(m + 1, 0));
    for (int i = 1; i <= n; i++)
      for (int j = 1; j <= m; j++) {
        if (a[i - 1] == b[j - 1])
          dp[i][j] = dp[i - 1][j - 1] + 1;
        else dp[i][j] = max(dp[i - 1][j], dp[i][j - 1]);
      }
  }
  int length() { return dp[n][m]; }
  string get() { // walk back from dp[n][m]
    string res;
    int i = n, j = m;
    while (i > 0 && j > 0) {
      if (a[i - 1] == b[j - 1]) res += a[--i], --j;
      else if (dp[i - 1][j] > dp[i][j - 1]) i--;
      else j--;
    }
    reverse(all(res));
    return res;
  }
};
\end{lstlisting}
\subsection{Edit distance}
\begin{lstlisting}
// Fewest insert / delete / replace operations that turn
// a into b. O(n * m)
int editDistance(const string &a, const string &b) {
  int n = a.size(), m = b.size();
  vector<vector<int>> dp(n + 1, vector<int>(m + 1));
  for (int i = 0; i <= n; i++) dp[i][0] = i; // delete
  for (int j = 0; j <= m; j++) dp[0][j] = j; // insert
  for (int i = 1; i <= n; i++)
    for (int j = 1; j <= m; j++)
      dp[i][j] = min({
        dp[i - 1][j] + 1, // delete a[i-1]
        dp[i][j - 1] + 1, // insert b[j-1]
        dp[i - 1][j - 1] + (a[i - 1] != b[j - 1])});
  return dp[n][m];
}
\end{lstlisting}
\subsection{Tree Rerooting}
\begin{lstlisting}
// Answer for EVERY node as the root in O(n).
// Example below: ans[u] = sum of distances from u to all
// nodes. CHANGE the two marked lines for other problems.
// Usage: fill adj and n, then dfsDown(root, -1) and
// dfsUp(root, -1).
const int N = 2e5 + 5;
vector<int> adj[N];
int n, sub[N], down[N], ans[N];
// down[u] = answer using only the subtree of u
void dfsDown(int u, int p) {
  sub[u] = 1, down[u] = 0;
  for (int v : adj[u]) {
    if (v == p) continue;
    dfsDown(v, u);
    sub[u] += sub[v];
    down[u] += down[v] + sub[v]; // CHANGE: merge child
  }
}
// ans[u] is known; move the root from u to child v
void dfsUp(int u, int p) {
  if (p == -1) ans[u] = down[u];
  for (int v : adj[u]) {
    if (v == p) continue;
    // CHANGE: sub[v] nodes get closer, the rest farther
    ans[v] = ans[u] - sub[v] + (n - sub[v]);
    dfsUp(v, u);
  }
}
\end{lstlisting}
\subsection{Convex Hull Trick (dynamic)}
\begin{lstlisting}
// Add lines y = k*x + m in any order, query the MAXIMUM
// at any x. O(log n) each. For MINIMUM: add(-k, -m) and
// use -query(x).
struct Line {
  mutable ll k, m, p; // p = last x where line is best
  bool operator<(const Line &o) const { return k < o.k; }
  bool operator<(ll x) const { return p < x; }
};
struct LineContainer : multiset<Line, less<>> {
  static const ll INF = LLONG_MAX;
  ll div(ll a, ll b) { // floored division
    return a / b - ((a ^ b) < 0 && a % b);
  }
  bool isect(iterator x, iterator y) {
    if (y == end()) return x->p = INF, 0;
    if (x->k == y->k) x->p = x->m > y->m ? INF : -INF;
    else x->p = div(y->m - x->m, x->k - y->k);
    return x->p >= y->p;
  }
  void add(ll k, ll m) {
    auto z = insert({k, m, 0}), y = z++, x = y;
    while (isect(y, z)) z = erase(z);
    if (x != begin() && isect(--x, y))
      isect(x, y = erase(y));
    while ((y = x) != begin() && (--x)->p >= y->p)
      isect(x, erase(y));
  }
  ll query(ll x) { // container must not be empty
    auto l = *lower_bound(x);
    return l.k * x + l.m;
  }
};
\end{lstlisting}
\subsection{Knapsack}
\begin{lstlisting}
// 0/1 knapsack: every item at most once. O(n * W).
// dp[w] = best value with total weight <= w
int knapsack01(const vector<int> &wt,
               const vector<int> &val, int W) {
  vector<int> dp(W + 1, 0);
  for (int i = 0; i < (int)wt.size(); i++)
    for (int w = W; w >= wt[i]; w--) // BACKWARD: once
      dp[w] = max(dp[w], dp[w - wt[i]] + val[i]);
  return dp[W];
}
// Unbounded (any number of each item): loop w FORWARD
//  for (int w = wt[i]; w <= W; w++) ...
// Bounded (item i at most c[i] times): split c[i] into
//  1, 2, 4, ..., rest and treat the parts as 0/1 items.
// Count ways (coin change), dp[0] = 1, dp[w] += dp[w-c]:
//  coins in the OUTER loop = combinations,
//  sum in the OUTER loop   = ordered ways.
// Subset sum with bitset, O(n * S / 64):
//  can[s] == 1 <=> some subset has sum s
const int MAXS = 1e5 + 5; // CHANGE: max total sum + 1
bitset<MAXS> subsetSums(const vector<int> &a) {
  bitset<MAXS> can;
  can[0] = 1;
  for (int x : a) can |= can << x;
  return can;
}
\end{lstlisting}
\subsection{Digit DP}
\begin{lstlisting}
// Count numbers x in [0, n] with a property of digits.
// Example below: digit sum == target.
// CHANGE the state "sum" and the check at the end.
string num;
int target;
int memoD[20][200][2]; // [position][sum][tight]
// pos = digit we fill now, tight = prefix is still equal
// to the prefix of n (so this digit is limited)
int digitDp(int pos, int sum, bool tight) {
  if (pos == (int)num.size()) return sum == target;
  int &res = memoD[pos][sum][tight];
  if (res != -1) return res;
  res = 0;
  int lim = tight ? num[pos] - '0' : 9;
  for (int d = 0; d <= lim; d++)
    res += digitDp(pos + 1, sum + d, tight && d == lim);
  return res;
}
int countUpTo(int n, int S) {
  if (n < 0) return 0;
  num = to_string(n), target = S;
  memset(memoD, -1, sizeof(memoD));
  return digitDp(0, 0, true);
}
// answer for [L, R]: countUpTo(R, S) - countUpTo(L-1, S)
// If leading zeros matter (e.g. counting digit 0), add a
// "started" flag to the state.
\end{lstlisting}
\subsection{Bitmask DP (TSP)}
\begin{lstlisting}
// Cheapest tour that starts at node 0, visits every node
// once and returns to 0. O(2^n * n^2), n <= 20.
// dp[mask][i] = min cost, visited set = mask, now at i
int tsp(const vector<vector<int>> &d) {
  int n = d.size(), full = (1LL << n) - 1;
  vector<vector<int>> dp(full + 1, vector<int>(n, inf));
  dp[1][0] = 0;
  for (int mask = 1; mask <= full; mask++)
    for (int i = 0; i < n; i++) {
      if (dp[mask][i] == inf) continue;
      for (int j = 0; j < n; j++) {
        if ((mask >> j) & 1) continue; // already visited
        int nm = mask | (1LL << j);
        dp[nm][j] =
          min(dp[nm][j], dp[mask][i] + d[i][j]);
      }
    }
  int best = inf;
  for (int i = 0; i < n; i++)
    best = min(best, dp[full][i] + d[i][0]);
  return best;
}
\end{lstlisting}
\subsection{SOS DP (sum over subsets)}
\begin{lstlisting}
// After the call f[mask] = sum of the old f[sub] over
// all submasks sub of mask. O(B * 2^B); f.size() == 2^B
void sos(vector<int> &f, int B) {
  for (int i = 0; i < B; i++)
    for (int mask = 0; mask < (1LL << B); mask++)
      if ((mask >> i) & 1)
        f[mask] += f[mask ^ (1LL << i)];
}
// Sum over SUPERSETS instead:
//  if (!((mask >> i) & 1))
//    f[mask] += f[mask | (1LL << i)];
// Undo (get the original back): same loops with -=
\end{lstlisting}
\subsection{Divide and conquer DP optimization}
\begin{lstlisting}
// dp[k][i] = min over j < i of dp[k-1][j] + cost(j+1, i)
// Works when the best j never moves left as i grows
// (true for most "split array into K parts" costs).
// O(K * n log n) instead of O(K * n^2). Items are 1..n.
vector<int> prv, cur; // dp[k-1][.] and dp[k][.]
int cost(int l, int r); // CHANGE: cost of group [l, r]
void compute(int l, int r, int optl, int optr) {
  if (l > r) return;
  int mid = (l + r) / 2;
  pair<int, int> best = {inf, optl}; // {value, best j}
  for (int j = optl; j <= min(mid - 1, optr); j++)
    if (prv[j] < inf)
      best = min(best, {prv[j] + cost(j + 1, mid), j});
  cur[mid] = best.first;
  compute(l, mid - 1, optl, best.second);
  compute(mid + 1, r, best.second, optr);
}
// min total cost of splitting items 1..n into K groups
int solveDC(int n, int K) {
  prv.assign(n + 1, inf), cur.assign(n + 1, inf);
  prv[0] = 0;
  for (int k = 1; k <= K; k++) {
    compute(1, n, 0, n - 1);
    cur[0] = inf;
    prv = cur;
  }
  return prv[n];
}
\end{lstlisting}
\section{Graph}
\subsection{DSU}
\begin{lstlisting}
// Works for 0- or 1-indexed nodes. ~O(1) per operation.
struct DSU {
  vector<int> parent, sz;
  int components;
  DSU(int n)
    : parent(n + 1), sz(n + 1, 1), components(n) {
    iota(all(parent), 0);
  }
  int find(int x) {
    if (x == parent[x]) return x;
    return parent[x] = find(parent[x]); // compress path
  }
  bool unite(int x, int y) { // false if already together
    x = find(x), y = find(y);
    if (x == y) return false;
    if (sz[x] < sz[y]) swap(x, y);
    parent[y] = x;
    sz[x] += sz[y];
    components--;
    return true;
  }
  bool same(int x, int y) { return find(x) == find(y); }
  int size(int x) { return sz[find(x)]; }
};
\end{lstlisting}
\subsection{DSU Rollback}
\begin{lstlisting}
// Undo the last unite(). O(log n) per operation (no path
// compression). Used in offline dynamic connectivity /
// divide and conquer over time.
// Usage: int t = d.snapshot(); ...unites...;
//  d.rollback(t);
struct DSURollback {
  vector<int> parent, sz;
  int components;
  vector<pair<int, int>> history; // {child, new root}
  DSURollback(int n)
    : parent(n + 1), sz(n + 1, 1), components(n) {
    iota(all(parent), 0);
  }
  int find(int x) {
    while (x != parent[x]) x = parent[x];
    return x;
  }
  bool unite(int x, int y) {
    x = find(x), y = find(y);
    if (x == y) return false;
    if (sz[x] < sz[y]) swap(x, y);
    history.push_back({y, x});
    parent[y] = x;
    sz[x] += sz[y];
    components--;
    return true;
  }
  int snapshot() { return history.size(); }
  // undo every unite made after snapshot() returned t
  void rollback(int t) {
    while ((int)history.size() > t) {
      auto [y, x] = history.back();
      history.pop_back();
      parent[y] = y;
      sz[x] -= sz[y];
      components++;
    }
  }
};
\end{lstlisting}
\subsection{Partially persistent DSU}
\begin{lstlisting}
// Ask "were u and v connected after the t-th unite?".
// Only the newest version can be changed. O(log n).
struct PersistentDSU {
  vector<int> parent, sz;
  vector<int> joined; // time when v got a parent
  int now = 0;
  PersistentDSU(int n)
    : parent(n + 1), sz(n + 1, 1), joined(n + 1, inf) {
    iota(all(parent), 0);
  }
  int find(int x, int t) { // root of x at time t
    while (x != parent[x] && joined[x] <= t)
      x = parent[x];
    return x;
  }
  void unite(int u, int v) { // unite number ++now
    now++;
    u = find(u, now), v = find(v, now);
    if (u == v) return;
    if (sz[u] < sz[v]) swap(u, v);
    parent[v] = u;
    sz[u] += sz[v];
    joined[v] = now;
  }
  bool same(int u, int v, int t) {
    return find(u, t) == find(v, t);
  }
};
\end{lstlisting}
\subsection{Fully persistent DSU}
\begin{lstlisting}
// Every unite() creates a NEW version; any old version
// can be queried or extended. O(log^2 n) per operation.
// Nodes are 0..n-1. Usage:
//  FullyPersistentDSU d(n); int v0 = d.init();
//  int v1 = d.unite(v0, a, b); d.find(v1, x);
struct FullyPersistentDSU {
  struct Node { int l, r, par, sz; };
  vector<Node> t; // persistent segment tree over nodes
  int n;
  FullyPersistentDSU(int n) : n(n) {}
  int build(int l, int r) {
    int id = t.size();
    t.push_back({-1, -1, l, 1});
    if (l == r) return id;
    int mid = (l + r) / 2;
    int a = build(l, mid), b = build(mid + 1, r);
    t[id].l = a, t[id].r = b;
    return id;
  }
  int init() { return build(0, n - 1); }
  // new version where node pos has parent par, size sz
  int update(int prev, int l, int r, int pos, int par,
             int sz) {
    int id = t.size();
    t.push_back(t[prev]);
    if (l == r) {
      t[id].par = par, t[id].sz = sz;
      return id;
    }
    int mid = (l + r) / 2, nw;
    if (pos <= mid) {
      nw = update(t[prev].l, l, mid, pos, par, sz);
      t[id].l = nw;
    } else {
      nw = update(t[prev].r, mid + 1, r, pos, par, sz);
      t[id].r = nw;
    }
    return id;
  }
  pair<int, int> query(int id, int l, int r, int pos) {
    while (l < r) {
      int mid = (l + r) / 2;
      if (pos <= mid) id = t[id].l, r = mid;
      else id = t[id].r, l = mid + 1;
    }
    return {t[id].par, t[id].sz};
  }
  int find(int ver, int u) {
    while (true) {
      int p = query(ver, 0, n - 1, u).first;
      if (p == u) return u;
      u = p;
    }
  }
  // returns the id of the new version
  int unite(int ver, int u, int v) {
    u = find(ver, u), v = find(ver, v);
    if (u == v) return ver;
    int su = query(ver, 0, n - 1, u).second;
    int sv = query(ver, 0, n - 1, v).second;
    if (su < sv) swap(u, v), swap(su, sv);
    ver = update(ver, 0, n - 1, v, u, sv);
    return update(ver, 0, n - 1, u, u, su + sv);
  }
};
\end{lstlisting}
\subsection{Dijkstra}
\begin{lstlisting}
// adj[u] = list of {v, w}, w >= 0. O((n + m) log n).
// Returns {dist, parent}. dist == inf: unreachable.
pair<vector<int>, vector<int>>
dijkstra(vector<vector<pair<int, int>>> &adj, int src) {
  int n = adj.size();
  vector<int> dist(n, inf), par(n, -1);
  priority_queue<pair<int, int>, vector<pair<int, int>>,
                 greater<pair<int, int>>> pq;
  // pq holds {distance, node}
  dist[src] = 0;
  pq.push({0, src});
  while (!pq.empty()) {
    auto [d, u] = pq.top();
    pq.pop();
    if (d > dist[u]) continue; // old entry, skip
    for (auto [v, w] : adj[u])
      if (dist[v] > d + w) {
        dist[v] = d + w;
        par[v] = u;
        pq.push({dist[v], v});
      }
  }
  return {dist, par};
}
// path to t: for (v = t; v != -1; v = par[v]), reverse
\end{lstlisting}
\subsection{Bellman Ford}
\begin{lstlisting}
// Negative edges allowed. O(n * m). Nodes 0..n or 1..n.
struct Edge { int u, v, w; };
struct BellmanFord {
  int n;
  vector<int> dist;
  bool negCycle = false; // reachable from the source
  BellmanFord(int n) : n(n), dist(n + 1, inf) {}
  void solve(int src, const vector<Edge> &edges) {
    dist[src] = 0;
    for (int i = 1; i <= n; i++) {
      bool changed = false;
      for (auto &e : edges)
        if (dist[e.u] != inf &&
            dist[e.u] + e.w < dist[e.v]) {
          dist[e.v] = dist[e.u] + e.w;
          changed = true;
        }
      if (!changed) return;
    }
    negCycle = true; // still improving after n rounds
  }
};
\end{lstlisting}
\subsection{Floyd Warshall}
\begin{lstlisting}
// All pairs shortest path. O(n^3). Nodes 1..n (or 0..n).
struct FloydWarshall {
  int n;
  vector<vector<int>> dist;
  FloydWarshall(int n)
    : n(n), dist(n + 1, vector<int>(n + 1, inf)) {
    for (int i = 0; i <= n; i++) dist[i][i] = 0;
  }
  void add_edge(int u, int v, int w, bool both = false) {
    dist[u][v] = min(dist[u][v], w);
    if (both) dist[v][u] = min(dist[v][u], w);
  }
  void solve() {
    for (int k = 0; k <= n; k++) // k = OUTER loop!
      for (int i = 0; i <= n; i++)
        for (int j = 0; j <= n; j++)
          if (dist[i][k] < inf && dist[k][j] < inf)
            dist[i][j] =
              min(dist[i][j], dist[i][k] + dist[k][j]);
  }
  bool negCycle() {
    for (int i = 0; i <= n; i++)
      if (dist[i][i] < 0) return true;
    return false;
  }
};
\end{lstlisting}
\subsection{Kruskal's MST}
\begin{lstlisting}
// edges = {u, v, w}. O(E log E). Needs DSU.
// Returns total weight; used = edges inside the MST.
// Graph is connected <=> used.size() == n - 1.
int kruskal(int n, vector<array<int, 3>> edges,
            vector<array<int, 3>> &used) {
  sort(all(edges), [](auto &a, auto &b) {
    return a[2] < b[2]; // cheapest edge first
  });
  DSU dsu(n);
  int total = 0;
  for (auto &e : edges)
    if (dsu.unite(e[0], e[1])) {
      total += e[2];
      used.push_back(e);
    }
  return total;
}
\end{lstlisting}
\subsection{Prim's MST}
\begin{lstlisting}
// adj[u] = list of {v, w}. O(E log V). Returns the total
// weight and fills mst with edges {parent, node, w}.
int prim(vector<vector<pair<int, int>>> &adj, int start,
         vector<array<int, 3>> &mst) {
  int n = adj.size(), total = 0;
  vector<bool> vis(n, false);
  priority_queue<array<int, 3>, vector<array<int, 3>>,
                 greater<array<int, 3>>> pq;
  // pq holds {weight, node, parent}
  pq.push({0, start, -1});
  while (!pq.empty()) {
    auto [w, u, par] = pq.top();
    pq.pop();
    if (vis[u]) continue;
    vis[u] = true;
    total += w;
    if (par != -1) mst.push_back({par, u, w});
    for (auto [v, wt] : adj[u])
      if (!vis[v]) pq.push({wt, v, u});
  }
  return total;
}
\end{lstlisting}
\subsection{Topological sort (Kahn)}
\begin{lstlisting}
// Nodes 1..n, adj = directed edges. O(n + m).
// Result has fewer than n nodes <=> there is a cycle.
// Smallest order: use a min priority_queue as the queue.
vector<int> topoSort(int n, vector<vector<int>> &adj) {
  vector<int> indeg(n + 1, 0), order;
  for (int u = 1; u <= n; u++)
    for (int v : adj[u]) indeg[v]++;
  queue<int> q;
  for (int u = 1; u <= n; u++)
    if (indeg[u] == 0) q.push(u);
  while (!q.empty()) {
    int u = q.front();
    q.pop();
    order.push_back(u);
    for (int v : adj[u])
      if (--indeg[v] == 0) q.push(v);
  }
  return order;
}
\end{lstlisting}
\subsection{SCC (Kosaraju) and condensation}
\begin{lstlisting}
// Directed graph, nodes 0..n-1. O(n + m).
// comp[v] = id of the component of v; ids are already in
// topological order (edge u->v: comp[u] <= comp[v]).
// dag = condensed graph (may hold repeated edges).
struct SCC {
  int n, cnt = 0;
  vector<vector<int>> adj, radj, dag;
  vector<int> comp, order;
  vector<bool> vis;
  SCC(int n)
    : n(n), adj(n), radj(n), comp(n, -1), vis(n) {}
  void add_edge(int u, int v) {
    adj[u].push_back(v);
    radj[v].push_back(u);
  }
  void dfs1(int v) { // order by exit time
    vis[v] = true;
    for (int u : adj[v])
      if (!vis[u]) dfs1(u);
    order.push_back(v);
  }
  void dfs2(int v) { // one component (reverse graph)
    comp[v] = cnt;
    for (int u : radj[v])
      if (comp[u] == -1) dfs2(u);
  }
  void build() {
    for (int i = 0; i < n; i++)
      if (!vis[i]) dfs1(i);
    reverse(all(order));
    for (int v : order)
      if (comp[v] == -1) dfs2(v), cnt++;
    dag.assign(cnt, {});
    for (int v = 0; v < n; v++)
      for (int u : adj[v])
        if (comp[v] != comp[u])
          dag[comp[v]].push_back(comp[u]);
  }
};
\end{lstlisting}
\subsection{Articulation points, bridges, bridge tree}
\begin{lstlisting}
// Undirected graph, multi-edges are handled (edge ids).
// Nodes 0..n or 1..n. O(n + m).
// Usage: CutBridge g(n); g.add_edge(u, v); g.run();
// then g.isCut[v], g.bridges; g.buildBridgeTree().
struct CutBridge {
  int n, m = 0, timer = 0;
  vector<vector<pair<int, int>>> adj; // {to, edge id}
  vector<int> tin, low; // low = highest node reachable
  vector<bool> isCut, isBridge; // isBridge[edge id]
  vector<pair<int, int>> bridges;
  CutBridge(int n)
    : n(n), adj(n + 1), tin(n + 1, -1), low(n + 1),
      isCut(n + 1, false) {}
  void add_edge(int u, int v) {
    adj[u].push_back({v, m});
    adj[v].push_back({u, m});
    m++;
  }
  void dfs(int v, int pe) { // pe = edge we came through
    tin[v] = low[v] = timer++;
    int children = 0;
    for (auto [to, id] : adj[v]) {
      if (id == pe) continue;
      if (tin[to] != -1) { // back edge
        low[v] = min(low[v], tin[to]);
        continue;
      }
      dfs(to, id);
      children++;
      low[v] = min(low[v], low[to]);
      if (low[to] > tin[v]) {
        isBridge[id] = true;
        bridges.push_back({v, to});
      }
      if (low[to] >= tin[v] && pe != -1) isCut[v] = true;
    }
    if (pe == -1 && children > 1) isCut[v] = true;
  }
  void run() {
    isBridge.assign(m, false);
    for (int i = 0; i <= n; i++)
      if (tin[i] == -1) dfs(i, -1);
  }
  // Bridge tree: comp[v] = 2-edge-connected component of
  // v, tree = graph of components joined by the bridges.
  // (an unused node 0 becomes its own component)
  vector<int> comp;
  vector<vector<int>> tree;
  void buildBridgeTree() {
    comp.assign(n + 1, -1);
    int c = 0;
    for (int i = 0; i <= n; i++) {
      if (comp[i] != -1) continue;
      vector<int> st = {i};
      comp[i] = c;
      while (!st.empty()) {
        int u = st.back();
        st.pop_back();
        for (auto [v, id] : adj[u])
          if (comp[v] == -1 && !isBridge[id])
            comp[v] = c, st.push_back(v);
      }
      c++;
    }
    tree.assign(c, {});
    for (auto [u, v] : bridges) {
      tree[comp[u]].push_back(comp[v]);
      tree[comp[v]].push_back(comp[u]);
    }
  }
};
\end{lstlisting}
\subsection{LCA (binary lifting)}
\begin{lstlisting}
// Build O(n log n), query O(log n). 0- or 1-indexed.
// Usage: LCA t(n, adj, root); t.lca(u, v); t.dist(u, v);
struct LCA {
  int n, LOG, timer = 0;
  vector<int> tin, tout, depth;
  vector<vector<int>> up; // up[v][i] = 2^i-th ancestor
  LCA(int n, vector<vector<int>> &adj, int root) : n(n) {
    LOG = 1;
    while ((1LL << LOG) <= n) LOG++;
    tin.assign(n + 1, 0), tout = depth = tin;
    up.assign(n + 1, vector<int>(LOG + 1));
    dfs(adj, root, root, 0);
  }
  void dfs(vector<vector<int>> &adj, int v, int p,
           int d) {
    tin[v] = ++timer;
    depth[v] = d;
    up[v][0] = p; // parent of the root = the root
    for (int i = 1; i <= LOG; ++i)
      up[v][i] = up[up[v][i - 1]][i - 1];
    for (int u : adj[v])
      if (u != p) dfs(adj, u, v, d + 1);
    tout[v] = ++timer;
  }
  // is u an ancestor of v (or u == v)?
  bool isAncestor(int u, int v) {
    return tin[u] <= tin[v] && tout[u] >= tout[v];
  }
  int lca(int u, int v) {
    if (isAncestor(u, v)) return u;
    if (isAncestor(v, u)) return v;
    for (int i = LOG; i >= 0; --i)
      if (!isAncestor(up[u][i], v)) u = up[u][i];
    return up[u][0];
  }
  int dist(int u, int v) { // number of edges on the path
    return depth[u] + depth[v] - 2 * depth[lca(u, v)];
  }
  int kthAncestor(int u, int k) { // k <= depth[u]
    for (int i = 0; i <= LOG; i++)
      if ((k >> i) & 1) u = up[u][i];
    return u;
  }
};
\end{lstlisting}
\subsection{Virtual tree}
\begin{lstlisting}
// Compress the tree to k chosen nodes plus their LCAs
// (at most 2k - 1 nodes). O(k log n) per call, so the
// sum over all queries stays small. Needs LCA.
// vadj must have size n + 1 (create it once, global).
// After the call "nodes" holds every node of the virtual
// tree; the returned node is its root.
vector<vector<int>> vadj;
int buildVirtualTree(vector<int> &nodes, LCA &t) {
  auto cmp = [&](int a, int b) {
    return t.tin[a] < t.tin[b];
  };
  sort(all(nodes), cmp);
  int k = nodes.size();
  for (int i = 0; i + 1 < k; i++) // LCA of neighbours
    nodes.push_back(t.lca(nodes[i], nodes[i + 1]));
  sort(all(nodes), cmp);
  nodes.erase(unique(all(nodes)), nodes.end());
  for (int u : nodes) vadj[u].clear();
  vector<int> stk = {nodes[0]}; // path from the root
  for (int i = 1; i < (int)nodes.size(); i++) {
    int u = nodes[i];
    while (!t.isAncestor(stk.back(), u)) stk.pop_back();
    vadj[stk.back()].push_back(u); // parent -> child
    stk.push_back(u);
  }
  return nodes[0];
}
// real length of an edge: depth[child] - depth[parent]
\end{lstlisting}
\subsection{HLD}
\begin{lstlisting}
// Path and subtree queries on a tree. O(log^2 n) each.
// Needs SegTree (Range Queries). Nodes 0..n-1.
// Usage: HLD h(n); h.add_edge(u, v) ...; h.init(root);
//  h.updatePath(u, v, x); h.queryPath(u, v);
// Values on EDGES: keep an edge's value at its deeper
// node and skip the top node: use pos[u] + 1 in the last
// update/query of updatePath/queryPath.
struct HLD {
  int n, curPos = 0;
  vector<vector<int>> adj;
  vector<int> parent, depth, heavy, head, pos, sz;
  SegTree st; // position pos[v] = value of node v
  HLD(int n)
    : n(n), adj(n), parent(n), depth(n), heavy(n, -1),
      head(n), pos(n), sz(n), st(n) {}
  void add_edge(int u, int v) {
    adj[u].push_back(v);
    adj[v].push_back(u);
  }
  void dfsSize(int v, int p, int d) {
    sz[v] = 1, parent[v] = p, depth[v] = d;
    int best = 0;
    for (int c : adj[v]) {
      if (c == p) continue;
      dfsSize(c, v, d + 1);
      sz[v] += sz[c];
      if (sz[c] > best) best = sz[c], heavy[v] = c;
    }
  }
  void dfsHld(int v, int h) { // h = top of the chain
    head[v] = h, pos[v] = curPos++;
    if (heavy[v] != -1) dfsHld(heavy[v], h);
    for (int c : adj[v])
      if (c != parent[v] && c != heavy[v]) dfsHld(c, c);
  }
  void init(int root = 0) {
    dfsSize(root, -1, 0);
    dfsHld(root, root);
  }
  void updatePath(int u, int v, int val) {
    for (; head[u] != head[v]; u = parent[head[u]]) {
      if (depth[head[u]] < depth[head[v]]) swap(u, v);
      st.update(pos[head[u]], pos[u], val);
    }
    if (depth[u] > depth[v]) swap(u, v);
    st.update(pos[u], pos[v], val);
  }
  int queryPath(int u, int v) {
    int res = 0;
    for (; head[u] != head[v]; u = parent[head[u]]) {
      if (depth[head[u]] < depth[head[v]]) swap(u, v);
      res += st.query(pos[head[u]], pos[u]);
    }
    if (depth[u] > depth[v]) swap(u, v);
    return res + st.query(pos[u], pos[v]);
  }
  // subtree of u = positions pos[u]..pos[u]+sz[u]-1
  void updateSubtree(int u, int val) {
    st.update(pos[u], pos[u] + sz[u] - 1, val);
  }
  int querySubtree(int u) {
    return st.query(pos[u], pos[u] + sz[u] - 1);
  }
  int lca(int u, int v) {
    for (; head[u] != head[v]; u = parent[head[u]])
      if (depth[head[u]] < depth[head[v]]) swap(u, v);
    return depth[u] < depth[v] ? u : v;
  }
};
\end{lstlisting}
\subsection{Dinic's max flow}
\begin{lstlisting}
// O(V^2 E) worst case; O(E sqrt(V)) on unit / bipartite
// graphs. Nodes 0..n-1.
// Usage: Dinic d(n, s, t); d.add_edge(u, v, cap);
//  d.maxFlow(); then d.minCut() / edge flows.
// Bipartite matching: s -> left (1), left -> right (1),
// right -> t (1); answer = maxFlow().
struct Dinic {
  struct FlowEdge { int v, u, cap, flow; };
  vector<FlowEdge> edges; // edge i and i^1 are a pair
  vector<vector<int>> adj;
  vector<int> level, ptr;
  int n, s, t;
  Dinic(int n, int s, int t)
    : adj(n), level(n), ptr(n), n(n), s(s), t(t) {}
  void add_edge(int v, int u, int cap) {
    adj[v].push_back(edges.size());
    edges.push_back({v, u, cap, 0});
    adj[u].push_back(edges.size());
    edges.push_back({u, v, 0, 0}); // reverse edge
    // undirected edge: give the reverse edge cap too
  }
  bool bfs() { // level graph from s
    fill(all(level), -1);
    level[s] = 0;
    queue<int> q;
    q.push(s);
    while (!q.empty()) {
      int v = q.front();
      q.pop();
      for (int id : adj[v]) {
        auto &e = edges[id];
        if (e.cap - e.flow == 0 || level[e.u] != -1)
          continue;
        level[e.u] = level[v] + 1;
        q.push(e.u);
      }
    }
    return level[t] != -1;
  }
  int dfs(int v, int pushed) {
    if (pushed == 0 || v == t) return pushed;
    for (int &i = ptr[v]; i < (int)adj[v].size(); i++) {
      int id = adj[v][i], u = edges[id].u;
      int left = edges[id].cap - edges[id].flow;
      if (level[v] + 1 != level[u] || !left) continue;
      int tr = dfs(u, min(pushed, left));
      if (tr == 0) continue;
      edges[id].flow += tr;
      edges[id ^ 1].flow -= tr;
      return tr;
    }
    return 0;
  }
  int maxFlow() {
    int f = 0;
    while (bfs()) {
      fill(all(ptr), 0);
      while (int pushed = dfs(s, inf)) f += pushed;
    }
    return f;
  }
  // call after maxFlow(): true = node is on the s side.
  // Cut edges: side[e.v] && !side[e.u] (even ids only)
  vector<bool> minCut() {
    vector<bool> side(n);
    for (int i = 0; i < n; i++) side[i] = level[i] != -1;
    return side;
  }
  // flow sent on the i-th added edge: edges[2 * i].flow
};
\end{lstlisting}
\subsection{Cartesian tree}
\begin{lstlisting}
// Min-heap on values, in-order = array order. O(n).
// The root is the position of the minimum; node i covers
// the maximal segment where a[i] is the minimum.
struct CartesianTree {
  int n, root = -1;
  vector<int> lc, rc; // children, -1 = none
  CartesianTree(const vector<int> &a)
    : n(a.size()), lc(n, -1), rc(n, -1) {
    vector<int> st; // increasing stack of indices
    for (int i = 0; i < n; i++) {
      int last = -1;
      while (!st.empty() && a[st.back()] > a[i]) {
        last = st.back();
        st.pop_back();
      }
      lc[i] = last;
      if (!st.empty()) rc[st.back()] = i;
      st.push_back(i);
    }
    if (n) root = st[0];
  }
};
\end{lstlisting}
\subsection{Grid BFS and 0-1 BFS}
\begin{lstlisting}
int dx[4] = {0, 1, 0, -1}, dy[4] = {1, 0, -1, 0};
// 8 directions: dx8 = {-1, -1, -1, 0, 0, 1, 1, 1},
//               dy8 = {-1, 0, 1, -1, 1, -1, 0, 1}
// Steps from (sr, sc) to every cell; '#' = wall; -1 = no
vector<vector<int>> gridBfs(const vector<string> &g,
                            int sr, int sc) {
  int n = g.size(), m = g[0].size();
  vector<vector<int>> dist(n, vector<int>(m, -1));
  queue<pair<int, int>> q;
  dist[sr][sc] = 0;
  q.push({sr, sc});
  while (!q.empty()) {
    auto [x, y] = q.front();
    q.pop();
    for (int d = 0; d < 4; d++) {
      int nx = x + dx[d], ny = y + dy[d];
      if (nx < 0 || nx >= n || ny < 0 || ny >= m)
        continue; // outside the grid
      if (g[nx][ny] == '#' || dist[nx][ny] != -1)
        continue; // wall or already visited
      dist[nx][ny] = dist[x][y] + 1;
      q.push({nx, ny});
    }
  }
  return dist;
}
// 0-1 BFS: every edge weight is 0 or 1. O(n + m).
// adj[u] = list of {v, w}
vector<int> bfs01(vector<vector<pair<int, int>>> &adj,
                  int src) {
  vector<int> dist(adj.size(), inf);
  deque<int> dq;
  dist[src] = 0;
  dq.push_back(src);
  while (!dq.empty()) {
    int u = dq.front();
    dq.pop_front();
    for (auto [v, w] : adj[u])
      if (dist[u] + w < dist[v]) {
        dist[v] = dist[u] + w;
        if (w == 0) dq.push_front(v); // same distance
        else dq.push_back(v);
      }
  }
  return dist;
}
\end{lstlisting}
\subsection{Tree diameter and Euler tour}
\begin{lstlisting}
// farthest node from src and all distances from src
pair<int, vector<int>> farthest(vector<vector<int>> &adj,
                                int src) {
  vector<int> dist(adj.size(), -1);
  queue<int> q;
  dist[src] = 0;
  q.push(src);
  int last = src;
  while (!q.empty()) {
    int u = last = q.front();
    q.pop();
    for (int v : adj[u])
      if (dist[v] == -1)
        dist[v] = dist[u] + 1, q.push(v);
  }
  return {last, dist};
}
// Longest path in a tree (number of edges). The two ends
// are a and b; farthest from any v = max(da[v], db[v])
int treeDiameter(vector<vector<int>> &adj, int anyNode) {
  int a = farthest(adj, anyNode).first;
  auto [b, da] = farthest(adj, a);
  return da[b];
}
// Euler tour: the subtree of v is exactly the positions
// tin[v]..tout[v], so a subtree query becomes a range
// query on a BIT / segment tree. Resize tin, tout first.
int timer = 0;
vector<int> tin, tout;
void eulerTour(vector<vector<int>> &adj, int v, int p) {
  tin[v] = timer++;
  for (int u : adj[v])
    if (u != p) eulerTour(adj, u, v);
  tout[v] = timer - 1;
}
\end{lstlisting}
\subsection{Small to large (DSU on tree)}
\begin{lstlisting}
// One answer per subtree in O(n log^2 n) total: always
// pour the SMALLER container into the bigger one.
// Example: number of different colours in every subtree.
const int N = 2e5 + 5;
vector<int> adj[N];
int col[N], res[N];
set<int> bag[N];
void dfsMerge(int u, int p) {
  bag[u].insert(col[u]);
  for (int v : adj[u]) {
    if (v == p) continue;
    dfsMerge(v, u);
    if (bag[v].size() > bag[u].size())
      swap(bag[u], bag[v]); // keep the big one in u
    for (int x : bag[v]) bag[u].insert(x);
    bag[v].clear();
  }
  res[u] = bag[u].size(); // CHANGE: answer of subtree u
}
\end{lstlisting}
\subsection{Centroid decomposition}
\begin{lstlisting}
// Remove the centroid, recurse on the pieces: every node
// has only O(log n) centroid ancestors. Used for path
// problems (count / best path): handle all paths that
// pass through each centroid c.
// cpar[v] = parent of v in the centroid tree.
// Usage: fill adj, then decompose(anyNode, -1).
const int N = 2e5 + 5;
vector<int> adj[N];
int sub[N], cpar[N];
bool removed[N];
int calcSize(int u, int p) {
  sub[u] = 1;
  for (int v : adj[u])
    if (v != p && !removed[v]) sub[u] += calcSize(v, u);
  return sub[u];
}
int findCentroid(int u, int p, int total) {
  for (int v : adj[u])
    if (v != p && !removed[v] && sub[v] * 2 > total)
      return findCentroid(v, u, total);
  return u;
}
void decompose(int u, int p) {
  int c = findCentroid(u, -1, calcSize(u, -1));
  cpar[c] = p;
  removed[c] = true;
  // CHANGE: c is the centroid of its piece. Do the work
  // here: dfs from c over the nodes with !removed[v].
  for (int v : adj[c])
    if (!removed[v]) decompose(v, c);
}
\end{lstlisting}
\subsection{2-SAT}
\begin{lstlisting}
// n boolean variables and clauses "(x is vx) OR (y is
// vy)". Finds an assignment or reports impossible.
// O(n + clauses). Needs SCC.
// Node 2*x means "x is true", 2*x+1 means "x is false".
struct TwoSat {
  int n;
  SCC scc;
  vector<bool> val; // the answer after solve()
  TwoSat(int n) : n(n), scc(2 * n) {}
  // at least one of (x == vx), (y == vy) must hold
  void addClause(int x, bool vx, int y, bool vy) {
    int a = 2 * x + !vx, b = 2 * y + !vy;
    scc.add_edge(a ^ 1, b); // if a is false, b must hold
    scc.add_edge(b ^ 1, a);
  }
  void force(int x, bool vx) { addClause(x, vx, x, vx); }
  // x -> y   : addClause(x, false, y, true)
  // x != y   : addClause(x, 1, y, 1) and (x, 0, y, 0)
  // x == y   : addClause(x, 1, y, 0) and (x, 0, y, 1)
  bool solve() {
    scc.build();
    val.assign(n, false);
    for (int x = 0; x < n; x++) {
      int t = scc.comp[2 * x], f = scc.comp[2 * x + 1];
      if (t == f) return false; // x and not x are equal
      val[x] = t > f; // later component = the true one
    }
    return true;
  }
};
\end{lstlisting}
\subsection{Bipartite matching (Hopcroft-Karp)}
\begin{lstlisting}
// Left nodes 0..n-1, right nodes 0..m-1. O(E sqrt(V)).
// Usage: HopcroftKarp hk(n, m); hk.add_edge(l, r);
//  hk.maxMatching(); matchL[l] = partner of l or -1.
// Konig: min vertex cover = max matching;
// max independent set = n + m - max matching;
// min path cover of a DAG = nodes - max matching.
struct HopcroftKarp {
  int n, m;
  vector<vector<int>> adj;
  vector<int> matchL, matchR, dist;
  HopcroftKarp(int n, int m)
    : n(n), m(m), adj(n), matchL(n, -1), matchR(m, -1),
      dist(n) {}
  void add_edge(int l, int r) { adj[l].push_back(r); }
  bool bfs() { // layers from all free left nodes
    queue<int> q;
    bool found = false;
    for (int l = 0; l < n; l++) {
      dist[l] = matchL[l] == -1 ? 0 : -1;
      if (dist[l] == 0) q.push(l);
    }
    while (!q.empty()) {
      int l = q.front();
      q.pop();
      for (int r : adj[l]) {
        int nl = matchR[r];
        if (nl == -1) found = true; // a free right node
        else if (dist[nl] == -1) {
          dist[nl] = dist[l] + 1;
          q.push(nl);
        }
      }
    }
    return found;
  }
  bool dfs(int l) {
    for (int r : adj[l]) {
      int nl = matchR[r];
      if (nl == -1 ||
          (dist[nl] == dist[l] + 1 && dfs(nl))) {
        matchL[l] = r, matchR[r] = l;
        return true;
      }
    }
    dist[l] = -1; // dead end: skip l in this round
    return false;
  }
  int maxMatching() {
    int res = 0;
    while (bfs())
      for (int l = 0; l < n; l++)
        if (matchL[l] == -1 && dfs(l)) res++;
    return res;
  }
};
\end{lstlisting}
\subsection{Min cost max flow}
\begin{lstlisting}
// Cheapest way to send the maximum flow. SPFA version:
// fine for V, E up to a few thousand. Negative costs are
// allowed (no negative cycle). Nodes 0..n-1.
// Usage: MCMF g(n); g.add_edge(u, v, cap, cost);
//  auto [flow, cost] = g.solve(s, t);
// Assignment problem: s -> worker (1, 0), worker -> job
// (1, price), job -> t (1, 0).
struct MCMF {
  struct E { int to, cap, cost; };
  int n;
  vector<E> edges; // edge i and i^1 are a pair
  vector<vector<int>> adj;
  vector<int> dist, par; // par[v] = edge used to reach v
  MCMF(int n) : n(n), adj(n) {}
  void add_edge(int u, int v, int cap, int cost) {
    adj[u].push_back(edges.size());
    edges.push_back({v, cap, cost});
    adj[v].push_back(edges.size());
    edges.push_back({u, 0, -cost});
  }
  bool spfa(int s, int t) { // cheapest path with room
    dist.assign(n, inf), par.assign(n, -1);
    vector<bool> inq(n, false);
    queue<int> q;
    dist[s] = 0;
    q.push(s);
    while (!q.empty()) {
      int u = q.front();
      q.pop();
      inq[u] = false;
      for (int id : adj[u]) {
        auto &e = edges[id];
        if (e.cap > 0 && dist[u] + e.cost < dist[e.to]) {
          dist[e.to] = dist[u] + e.cost;
          par[e.to] = id;
          if (!inq[e.to]) inq[e.to] = true, q.push(e.to);
        }
      }
    }
    return dist[t] != inf;
  }
  // sends at most maxFlow units, returns {flow, cost}
  pair<int, int> solve(int s, int t, int maxFlow = inf) {
    int flow = 0, cost = 0;
    while (flow < maxFlow && spfa(s, t)) {
      int f = maxFlow - flow; // bottleneck of the path
      for (int v = t; v != s; v = edges[par[v] ^ 1].to)
        f = min(f, edges[par[v]].cap);
      for (int v = t; v != s; v = edges[par[v] ^ 1].to) {
        edges[par[v]].cap -= f;
        edges[par[v] ^ 1].cap += f;
      }
      flow += f;
      cost += f * dist[t];
    }
    return {flow, cost};
  }
};
\end{lstlisting}
\subsection{Euler path / circuit (Hierholzer)}
\begin{lstlisting}
// A walk that uses EVERY edge exactly once. O(n + m).
// Undirected graph: exists <=> the edges are connected
// and 0 odd-degree nodes (circuit) or 2 (path between
// them). Returns the m + 1 nodes of the walk, or {}.
// Directed: need in == out everywhere (circuit), or one
// node with out = in + 1 (start), one with in = out + 1.
vector<int> eulerPath(int n,
                      vector<pair<int, int>> &edges) {
  int m = edges.size();
  vector<vector<pair<int, int>>> adj(n); // {to, edge id}
  for (int i = 0; i < m; i++) {
    adj[edges[i].first].push_back({edges[i].second, i});
    adj[edges[i].second].push_back({edges[i].first, i});
  }
  int start = m ? edges[0].first : 0, odd = 0;
  for (int v = 0; v < n; v++)
    if (adj[v].size() % 2) odd++, start = v;
  if (odd != 0 && odd != 2) return {};
  vector<bool> used(m, false);
  vector<int> ptr(n, 0), st = {start}, path;
  while (!st.empty()) {
    int u = st.back();
    while (ptr[u] < (int)adj[u].size() &&
           used[adj[u][ptr[u]].second])
      ptr[u]++; // skip edges that were already walked
    if (ptr[u] == (int)adj[u].size()) {
      path.push_back(u); // stuck: u is final in the walk
      st.pop_back();
    } else {
      auto [v, id] = adj[u][ptr[u]];
      used[id] = true;
      st.push_back(v);
    }
  }
  if ((int)path.size() != m + 1) return {}; // not joined
  reverse(all(path));
  return path;
}
\end{lstlisting}
\section{Range Queries}
\subsection{Prefix sum 2D}
\begin{lstlisting}
// Build O(n*m), query O(1). Cells are 0-based.
struct Pref2D {
  int n, m;
  vector<vector<int>> p; // p[i][j] = sum a[0..i)[0..j)
  Pref2D(const vector<vector<int>> &a)
    : n(a.size()), m(a[0].size()),
      p(n + 1, vector<int>(m + 1, 0)) {
    for (int i = 1; i <= n; ++i)
      for (int j = 1; j <= m; ++j)
        p[i][j] = a[i - 1][j - 1] + p[i - 1][j] +
                  p[i][j - 1] - p[i - 1][j - 1];
  }
  // sum of rows r1..r2, columns c1..c2 (inclusive)
  int sum(int r1, int c1, int r2, int c2) {
    if (r1 > r2 || c1 > c2) return 0;
    return p[r2 + 1][c2 + 1] - p[r1][c2 + 1] -
           p[r2 + 1][c1] + p[r1][c1];
  }
};
// 2D difference array (add v on a rectangle, read last):
//  d[r1][c1] += v; d[r1][c2+1] -= v; d[r2+1][c1] -= v;
//  d[r2+1][c2+1] += v;  then take 2D prefix sums of d.
\end{lstlisting}
\subsection{Sparse Table}
\begin{lstlisting}
// Static array, build O(n log n), query O(1) for min /
// max / gcd (idempotent). CHANGE merge() for max or gcd.
struct SparseTable {
  int n;
  vector<vector<int>> st; // st[j][i] covers [i, i+2^j)
  int merge(int a, int b) { return min(a, b); }
  SparseTable(const vector<int> &a) : n(a.size()) {
    st.assign(1, a);
    for (int j = 1; (1LL << j) <= n; j++) {
      st.push_back(vector<int>(n));
      for (int i = 0; i + (1LL << j) <= n; i++)
        st[j][i] = merge(st[j - 1][i],
          st[j - 1][i + (1LL << (j - 1))]);
    }
  }
  int query(int l, int r) { // [l, r] inclusive, l <= r
    int j = 63 - __builtin_clzll(r - l + 1);
    return merge(st[j][l], st[j][r - (1LL << j) + 1]);
  }
};
\end{lstlisting}
\subsection{Fenwick tree (BIT)}
\begin{lstlisting}
// Point add, range sum. O(log n). Indices are 0-based.
struct Fenwick {
  int n;
  vector<int> t;
  Fenwick(int n) : n(n), t(n + 1, 0) {}
  void add(int i, int delta) { // a[i] += delta
    for (++i; i <= n; i += i & -i) t[i] += delta;
  }
  int query(int i) { // sum of a[0..i]
    int sum = 0;
    for (++i; i > 0; i -= i & -i) sum += t[i];
    return sum;
  }
  int query(int l, int r) { // sum of a[l..r]
    return l > r ? 0 : query(r) - query(l - 1);
  }
  // smallest index i with a[0] + ... + a[i] >= target
  // (returns n if the total is smaller). Needs a[i]>=0
  int lowerBound(int target) {
    int pos = 0, pw = 1;
    while (pw * 2 <= n) pw *= 2;
    for (; pw > 0; pw >>= 1)
      if (pos + pw <= n && t[pos + pw] < target) {
        pos += pw;
        target -= t[pos];
      }
    return pos;
  }
};
// Range add + point query: add(l, v), add(r + 1, -v);
// the value of a[i] is then query(i).
// Range add + range sum with two BITs. Indices 0-based.
struct FenwickRange {
  Fenwick b1, b2;
  FenwickRange(int n) : b1(n + 1), b2(n + 1) {}
  void add(int l, int r, int v) { // a[l..r] += v
    b1.add(l, v), b1.add(r + 1, -v);
    b2.add(l, v * l), b2.add(r + 1, -v * (r + 1));
  }
  int prefix(int i) { // sum of a[0..i]
    return b1.query(i) * (i + 1) - b2.query(i);
  }
  int query(int l, int r) { // sum of a[l..r]
    return prefix(r) - (l ? prefix(l - 1) : 0);
  }
};
\end{lstlisting}
\subsection{Inversion count}
\begin{lstlisting}
// Pairs i < j with a[i] > a[j]. O(n log n).
// Needs Fenwick.
int inversionCount(vector<int> a) {
  vector<int> srt = a; // compress values to 0..k-1
  sort(all(srt));
  srt.erase(unique(all(srt)), srt.end());
  Fenwick fw(srt.size());
  int res = 0;
  for (int i = (int)a.size() - 1; i >= 0; i--) {
    int x = lower_bound(all(srt), a[i]) - srt.begin();
    res += fw.query(x - 1); // smaller, on the right
    fw.add(x, 1);
  }
  return res;
}
\end{lstlisting}
\subsection{2D BIT (range add, range sum)}
\begin{lstlisting}
// All operations O(log n * log m). Cells are 1-based.
struct BIT2D {
  int n, m;
  vector<vector<int>> t1, t2, t3, t4;
  BIT2D(int n, int m) : n(n), m(m) {
    t1.assign(n + 2, vector<int>(m + 2, 0));
    t2 = t3 = t4 = t1;
  }
  void upd(int r, int c, int v) {
    for (int i = r; i <= n; i += i & -i)
      for (int j = c; j <= m; j += j & -j) {
        t1[i][j] += v;
        t2[i][j] += v * (r - 1);
        t3[i][j] += v * (c - 1);
        t4[i][j] += v * (r - 1) * (c - 1);
      }
  }
  // add v to every cell of rectangle (r1,c1)..(r2,c2)
  void rangeAdd(int r1, int c1, int r2, int c2, int v) {
    upd(r1, c1, v), upd(r1, c2 + 1, -v);
    upd(r2 + 1, c1, -v), upd(r2 + 1, c2 + 1, v);
  }
  int pref(int r, int c) { // sum of (1,1)..(r,c)
    int res = 0;
    for (int i = r; i > 0; i -= i & -i)
      for (int j = c; j > 0; j -= j & -j)
        res += t1[i][j] * r * c - t2[i][j] * c -
               t3[i][j] * r + t4[i][j];
    return res;
  }
  int rangeSum(int r1, int c1, int r2, int c2) {
    return pref(r2, c2) - pref(r1 - 1, c2) -
           pref(r2, c1 - 1) + pref(r1 - 1, c1 - 1);
  }
};
// Only point add + rectangle sum: keep t1 alone.
\end{lstlisting}
\subsection{Segment tree (lazy)}
\begin{lstlisting}
// Range ADD, range SUM. O(log n). Indices 0..n-1.
// Usage: SegTree st(n); st.build(a);
//  st.update(l, r, x); st.query(l, r);
// CHANGE the 4 marked places for min / max / assign.
struct SegTree {
  int n;
  vector<int> t, lazy;
  SegTree(int n) : n(n), t(4 * n + 5), lazy(4 * n + 5) {}
  // (1) how two children combine; for max: max(a, b)
  int merge(int a, int b) { return a + b; }
  // (2) value returned for an empty range; for max: -inf
  int identity() { return 0; }
  // (3) add x to every cell of a node with len cells;
  //     for max/min: t[v] += x (no len). For ASSIGN:
  //     t[v] = x * len, and use a "nothing" marker
  //     (not 0) for an empty lazy
  void apply(int v, int len, int x) {
    t[v] += x * len;
    lazy[v] += x; // (4) for ASSIGN: lazy[v] = x
  }
  void push(int v, int tl, int tr) {
    if (lazy[v] == 0) return;
    int tm = (tl + tr) / 2;
    apply(2 * v, tm - tl + 1, lazy[v]);
    apply(2 * v + 1, tr - tm, lazy[v]);
    lazy[v] = 0;
  }
  void build(const vector<int> &a, int v, int tl,
             int tr) {
    if (tl == tr) { t[v] = a[tl]; return; }
    int tm = (tl + tr) / 2;
    build(a, 2 * v, tl, tm);
    build(a, 2 * v + 1, tm + 1, tr);
    t[v] = merge(t[2 * v], t[2 * v + 1]);
  }
  void update(int v, int tl, int tr, int l, int r,
              int x) {
    if (l > r) return;
    if (l == tl && r == tr) {
      apply(v, tr - tl + 1, x);
      return;
    }
    push(v, tl, tr);
    int tm = (tl + tr) / 2;
    update(2 * v, tl, tm, l, min(r, tm), x);
    update(2 * v + 1, tm + 1, tr, max(l, tm + 1), r, x);
    t[v] = merge(t[2 * v], t[2 * v + 1]);
  }
  int query(int v, int tl, int tr, int l, int r) {
    if (l > r) return identity();
    if (l == tl && r == tr) return t[v];
    push(v, tl, tr);
    int tm = (tl + tr) / 2;
    return merge(
      query(2 * v, tl, tm, l, min(r, tm)),
      query(2 * v + 1, tm + 1, tr, max(l, tm + 1), r));
  }
  // short forms that you call from outside
  void build(const vector<int> &a) {
    build(a, 1, 0, n - 1);
  }
  void update(int l, int r, int x) {
    update(1, 0, n - 1, l, r, x);
  }
  int query(int l, int r) {
    return query(1, 0, n - 1, l, r);
  }
};
\end{lstlisting}
\subsection{Segment tree (iterative)}
\begin{lstlisting}
// Point update, range query, no recursion: about 3x
// faster than the lazy tree. Query range is [l, r).
struct SegTreeIter {
  int n;
  vector<int> t;
  int merge(int a, int b) { return a + b; } // CHANGE
  int identity = 0; // CHANGE: inf for min, -inf for max
  SegTreeIter(int n) : n(n), t(2 * n, 0) {}
  void build(const vector<int> &a) {
    for (int i = 0; i < n; ++i) t[n + i] = a[i];
    for (int i = n - 1; i > 0; --i)
      t[i] = merge(t[2 * i], t[2 * i + 1]);
  }
  void set(int p, int value) { // a[p] = value
    for (t[p += n] = value; p > 1; p >>= 1)
      t[p >> 1] = merge(t[p & ~1LL], t[p | 1]);
  }
  int query(int l, int r) { // merge of a[l..r-1]
    int resl = identity, resr = identity;
    for (l += n, r += n; l < r; l >>= 1, r >>= 1) {
      if (l & 1) resl = merge(resl, t[l++]);
      if (r & 1) resr = merge(t[--r], resr);
    }
    return merge(resl, resr);
  }
};
\end{lstlisting}
\subsection{Dynamic segment tree}
\begin{lstlisting}
// Range add, range sum over a HUGE index range (1e9+)
// without compression: nodes are created when touched.
// ~2 * log(n) new nodes per update. merge() joins two
// trees (small-to-large on trees / sets of values).
// Usage: DynSegTree st(1e9); st.update(l, r, x);
struct DynSegTree {
  struct Node { int sum = 0, lazy = 0, l = 0, r = 0; };
  vector<Node> t; // t[0] = empty node, t[1] = root
  int n;
  DynSegTree(int n, int reserveNodes = 0) : t(2), n(n) {
    t.reserve(reserveNodes); // optional, avoids realloc
  }
  int newNode() {
    t.push_back(Node());
    return t.size() - 1;
  }
  void apply(int v, int len, int x) {
    t[v].sum += x * len;
    t[v].lazy += x;
  }
  void push(int v, int tl, int tr) {
    if (!t[v].l) { int c = newNode(); t[v].l = c; }
    if (!t[v].r) { int c = newNode(); t[v].r = c; }
    int tm = tl + (tr - tl) / 2;
    apply(t[v].l, tm - tl + 1, t[v].lazy);
    apply(t[v].r, tr - tm, t[v].lazy);
    t[v].lazy = 0;
  }
  void update(int v, int tl, int tr, int l, int r,
              int x) {
    if (l > r) return;
    if (l == tl && r == tr) {
      apply(v, tr - tl + 1, x);
      return;
    }
    push(v, tl, tr);
    int tm = tl + (tr - tl) / 2;
    update(t[v].l, tl, tm, l, min(r, tm), x);
    update(t[v].r, tm + 1, tr, max(l, tm + 1), r, x);
    t[v].sum = t[t[v].l].sum + t[t[v].r].sum;
  }
  int query(int v, int tl, int tr, int l, int r) {
    if (l > r || v == 0) return 0;
    if (l == tl && r == tr) return t[v].sum;
    push(v, tl, tr);
    int tm = tl + (tr - tl) / 2;
    return query(t[v].l, tl, tm, l, min(r, tm)) +
           query(t[v].r, tm + 1, tr, max(l, tm + 1), r);
  }
  void update(int l, int r, int x) {
    update(1, 0, n - 1, l, r, x);
  }
  int query(int l, int r) {
    return query(1, 0, n - 1, l, r);
  }
  // Merge tree v into tree u (roots made by newNode(),
  // POINT updates only, no lazy). Returns the new root.
  int merge(int u, int v, int tl, int tr) {
    if (!u || !v) return u ? u : v;
    if (tl == tr) {
      t[u].sum += t[v].sum;
      return u;
    }
    int tm = tl + (tr - tl) / 2;
    int a = merge(t[u].l, t[v].l, tl, tm);
    int b = merge(t[u].r, t[v].r, tm + 1, tr);
    t[u].l = a, t[u].r = b;
    t[u].sum = t[a].sum + t[b].sum;
    return u;
  }
};
\end{lstlisting}
\subsection{Persistent segment tree}
\begin{lstlisting}
// Every update makes a new version (root) in O(log n)
// and old versions stay usable. Point add, range sum.
// Memory: ~log(n) nodes per update.
// Usage: root[0] = 0 (empty tree);
//  root[i] = ps.update(root[i-1], 0, MX, pos, +1);
//  ps.query(root[i], 0, MX, l, r);
// k-th smallest in a[l..r]: build versions over VALUES,
//  ps.kth(root[l], root[r + 1], 0, MX, k)
struct PersistentSeg {
  struct Node { int sum = 0, l = 0, r = 0; };
  vector<Node> t; // t[0] = empty tree
  PersistentSeg() : t(1) {}
  int update(int prev, int tl, int tr, int pos,
             int val) {
    int cur = t.size();
    t.push_back(t[prev]); // copy the old node
    t[cur].sum += val;
    if (tl == tr) return cur;
    int tm = (tl + tr) / 2;
    if (pos <= tm) {
      int c = update(t[prev].l, tl, tm, pos, val);
      t[cur].l = c;
    } else {
      int c = update(t[prev].r, tm + 1, tr, pos, val);
      t[cur].r = c;
    }
    return cur;
  }
  int query(int cur, int tl, int tr, int l, int r) {
    if (!cur || l > r) return 0;
    if (l == tl && r == tr) return t[cur].sum;
    int tm = (tl + tr) / 2;
    return query(t[cur].l, tl, tm, l, min(r, tm)) +
           query(t[cur].r, tm + 1, tr, max(l, tm + 1),
                 r);
  }
  // k-th smallest (k >= 1) among values inserted between
  // version a (exclusive) and version b (inclusive)
  int kth(int a, int b, int tl, int tr, int k) {
    if (tl == tr) return tl;
    int tm = (tl + tr) / 2;
    int left = t[t[b].l].sum - t[t[a].l].sum;
    if (k <= left) return kth(t[a].l, t[b].l, tl, tm, k);
    return kth(t[a].r, t[b].r, tm + 1, tr, k - left);
  }
};
\end{lstlisting}
\subsection{Li Chao tree}
\begin{lstlisting}
// Add lines y = m*x + c, query the MINIMUM y at integer
// x in [lo, hi]. O(log(hi - lo)). For MAXIMUM: add
// (-m, -c) and negate the answer.
// Usage: LiChao lc(0, 1e9); lc.add({m, c}); lc.query(x);
struct LiChao {
  struct Line {
    int m = 0, c = inf;
    int at(int x) { return m * x + c; }
  };
  struct Node { Line line; int l = 0, r = 0; };
  vector<Node> t; // t[0] = empty, t[1] = root
  int lo, hi;
  LiChao(int lo, int hi) : t(2), lo(lo), hi(hi) {}
  int newNode() {
    t.push_back({});
    return t.size() - 1;
  }
  void add(int v, int l, int r, Line nw) {
    int mid = l + (r - l) / 2;
    // keep in this node the line that is better at mid
    if (nw.at(mid) < t[v].line.at(mid))
      swap(nw, t[v].line);
    if (l == r) return;
    // the loser can still win on ONE side only
    if (nw.at(l) < t[v].line.at(l)) {
      int c = t[v].l ? t[v].l : newNode();
      t[v].l = c;
      add(c, l, mid, nw);
    } else if (nw.at(r) < t[v].line.at(r)) {
      int c = t[v].r ? t[v].r : newNode();
      t[v].r = c;
      add(c, mid + 1, r, nw);
    }
  }
  void add(Line nw) { add(1, lo, hi, nw); }
  int query(int x) {
    int res = inf, v = 1, l = lo, r = hi;
    while (v) {
      res = min(res, t[v].line.at(x));
      int mid = l + (r - l) / 2;
      if (x <= mid) v = t[v].l, r = mid;
      else v = t[v].r, l = mid + 1;
    }
    return res;
  }
};
\end{lstlisting}
\subsection{Mo's Algorithm}
\begin{lstlisting}
// Offline range queries in O((n + q) * sqrt(n)).
// CHANGE add() / remove() so that "cur" is always the
// answer for the current window [curL, curR].
// Usage: Mo mo(a); mo.addQuery(l, r) ...; mo.solve();
//  then mo.ans[i] is the answer of the i-th query.
struct Mo {
  struct Query { int l, r, id; };
  int n, block;
  vector<int> a, ans;
  vector<Query> qs;
  int cur = 0;
  vector<int> freq; // example: number of distinct values
  Mo(const vector<int> &arr) : n(arr.size()), a(arr) {
    block = max<int>(1, sqrt(n));
    freq.assign(*max_element(all(a)) + 1, 0);
  }
  void addQuery(int l, int r) { // 0-based, inclusive
    qs.push_back({l, r, (int)qs.size()});
  }
  void add(int i) { // a[i] enters the window
    if (++freq[a[i]] == 1) cur++;
  }
  void remove(int i) { // a[i] leaves the window
    if (--freq[a[i]] == 0) cur--;
  }
  void solve() {
    sort(all(qs), [&](const Query &x, const Query &y) {
      int bx = x.l / block, by = y.l / block;
      if (bx != by) return bx < by;
      return (bx & 1) ? x.r > y.r : x.r < y.r; // zig-zag
    });
    ans.assign(qs.size(), 0);
    int curL = 0, curR = -1;
    for (auto &q : qs) {
      while (curL > q.l) add(--curL);
      while (curR < q.r) add(++curR);
      while (curL < q.l) remove(curL++);
      while (curR > q.r) remove(curR--);
      ans[q.id] = cur;
    }
  }
};
\end{lstlisting}
\subsection{Sqrt decomposition}
\begin{lstlisting}
// Point assign, range MIN. O(sqrt(n)) per operation.
// CHANGE min / identity for sum, max, ...
struct SqrtDecomp {
  int n, bs; // bs = block size
  vector<int> a, blk; // blk[b] = min of block b
  const int identity = inf;
  SqrtDecomp(const vector<int> &arr)
    : n(arr.size()), a(arr) {
    bs = sqrt(n) + 1;
    blk.assign(n / bs + 1, identity);
    for (int i = 0; i < n; ++i)
      blk[i / bs] = min(blk[i / bs], a[i]);
  }
  void update(int idx, int val) { // a[idx] = val
    int b = idx / bs;
    a[idx] = val;
    blk[b] = identity; // rebuild this block
    for (int i = b * bs; i < min(n, (b + 1) * bs); ++i)
      blk[b] = min(blk[b], a[i]);
  }
  int query(int l, int r) { // min of a[l..r]
    int res = identity, bl = l / bs, br = r / bs;
    if (bl == br) {
      for (int i = l; i <= r; ++i) res = min(res, a[i]);
      return res;
    }
    for (int i = l; i < (bl + 1) * bs; ++i)
      res = min(res, a[i]);        // left partial block
    for (int b = bl + 1; b < br; ++b)
      res = min(res, blk[b]);      // full blocks
    for (int i = br * bs; i <= r; ++i)
      res = min(res, a[i]);        // right partial block
    return res;
  }
};
\end{lstlisting}
\section{Geometry}
\subsection{Point and vector basics}
\begin{lstlisting}
const double EPS = 1e-9;
const double PI = acos(-1.0);
// compare doubles: -1 (a < b), 0 (equal), 1 (a > b)
int dcmp(double a, double b) {
  if (fabs(a - b) < EPS) return 0;
  return a < b ? -1 : 1;
}
double toRad(double deg) { return deg * PI / 180.0; }
double toDeg(double rad) { return rad * 180.0 / PI; }
struct Point {
  double x, y;
  Point(double x = 0, double y = 0) : x(x), y(y) {}
  Point operator+(const Point &o) const {
    return {x + o.x, y + o.y};
  }
  Point operator-(const Point &o) const {
    return {x - o.x, y - o.y};
  }
  Point operator*(double s) const {
    return {x * s, y * s};
  }
  Point operator/(double s) const {
    return {x / s, y / s};
  }
  bool operator<(const Point &o) const { // by x, then y
    if (dcmp(x, o.x) != 0) return x < o.x;
    return dcmp(y, o.y) < 0;
  }
  bool operator==(const Point &o) const {
    return dcmp(x, o.x) == 0 && dcmp(y, o.y) == 0;
  }
};
// dot > 0: angle < 90, dot == 0: perpendicular
double dot(Point a, Point b) {
  return a.x * b.x + a.y * b.y;
}
// cross > 0: b is counter-clockwise (left) of a;
// |cross| = area of the parallelogram
double cross(Point a, Point b) {
  return a.x * b.y - a.y * b.x;
}
double magSq(Point a) { return dot(a, a); }
double len(Point a) { return sqrt(dot(a, a)); }
double dist(Point a, Point b) { return len(a - b); }
// angle of the vector in (-PI, PI]
double angle(Point a) { return atan2(a.y, a.x); }
// p rotated by 90 degrees CCW
Point perp(Point p) { return {-p.y, p.x}; }
Point normalize(Point v) { // unit vector
  double l = len(v);
  return dcmp(l, 0) == 0 ? Point(0, 0) : v / l;
}
// rotate p by theta radians CCW around the origin
Point rotate(Point p, double theta) {
  return {p.x * cos(theta) - p.y * sin(theta),
          p.x * sin(theta) + p.y * cos(theta)};
}
Point rotateAround(Point p, double theta, Point pivot) {
  return pivot + rotate(p - pivot, theta);
}
// turn a -> b -> c: 0 = collinear, 1 = clockwise
// (right turn), 2 = counter-clockwise (left turn)
int orientation(Point a, Point b, Point c) {
  double val = cross(b - a, c - a);
  if (fabs(val) < EPS) return 0;
  return val < 0 ? 1 : 2;
}
// direction of the bisector of angle a-o-b
Point angleBisector(Point a, Point o, Point b) {
  return normalize(a - o) + normalize(b - o);
}
\end{lstlisting}
\subsection{Lines, segments and rays}
\begin{lstlisting}
// Needs "Point and vector basics". A line / segment /
// ray is given by two points a, b (a ray starts at a).
Point projectOnLine(Point p, Point a, Point b) {
  Point ab = b - a;
  return a + ab * (dot(p - a, ab) / dot(ab, ab));
}
Point reflectOverLine(Point p, Point a, Point b) {
  return projectOnLine(p, a, b) * 2 - p;
}
double distPointToLine(Point p, Point a, Point b) {
  return fabs(cross(b - a, p - a)) / dist(a, b);
}
double distPointToSegment(Point p, Point a, Point b) {
  double l2 = magSq(b - a);
  if (l2 < EPS) return dist(p, a); // a == b
  double t = dot(p - a, b - a) / l2;
  if (t < 0.0) return dist(p, a); // before a
  if (t > 1.0) return dist(p, b); // after b
  return distPointToLine(p, a, b);
}
double distPointToRay(Point p, Point a, Point b) {
  if (dot(p - a, b - a) < 0.0) return dist(p, a);
  return distPointToLine(p, a, b);
}
// is q on the segment p-r (end points included)?
bool onSegment(Point p, Point q, Point r) {
  return orientation(p, q, r) == 0 &&
         q.x >= min(p.x, r.x) - EPS &&
         q.x <= max(p.x, r.x) + EPS &&
         q.y >= min(p.y, r.y) - EPS &&
         q.y <= max(p.y, r.y) + EPS;
}
// LINES ab and cd: 0 = parallel, 1 = same line,
// 2 = one intersection point (stored in res)
int lineIntersection(Point a, Point b, Point c, Point d,
                     Point &res) {
  Point v1 = b - a, v2 = d - c;
  double det = cross(v1, v2);
  if (dcmp(det, 0) == 0)
    return dcmp(cross(v1, c - a), 0) == 0 ? 1 : 0;
  res = a + v1 * (cross(c - a, v2) / det);
  return 2;
}
bool segmentsIntersect(Point a, Point b, Point c,
                       Point d) {
  int o1 = orientation(a, b, c);
  int o2 = orientation(a, b, d);
  int o3 = orientation(c, d, a);
  int o4 = orientation(c, d, b);
  if (o1 != o2 && o3 != o4) return true; // they cross
  if (o1 == 0 && onSegment(a, c, b)) return true;
  if (o2 == 0 && onSegment(a, d, b)) return true;
  if (o3 == 0 && onSegment(c, a, d)) return true;
  if (o4 == 0 && onSegment(c, b, d)) return true;
  return false;
}
// intersection point of SEGMENTS ab and cd. Empty if
// they do not meet; if they overlap, one end of the
// common part.
vector<Point> segmentIntersection(Point a, Point b,
                                  Point c, Point d) {
  vector<Point> res;
  Point ab = b - a, cd = d - c, ac = c - a;
  double det = cross(ab, cd);
  if (fabs(det) < EPS) { // parallel
    if (fabs(cross(ab, ac)) > EPS ||
        fabs(cross(cd, ac)) > EPS)
      return res; // not on the same line
    Point s = max(min(a, b), min(c, d));
    Point e = min(max(a, b), max(c, d));
    if (s < e || s == e) res.push_back(e);
    return res;
  }
  double t = cross(ac, cd) / det; // position on ab
  double u = cross(ac, ab) / det; // position on cd
  if (t >= -EPS && t <= 1 + EPS && u >= -EPS &&
      u <= 1 + EPS)
    res.push_back(a + ab * min(1.0, max(0.0, t)));
  return res;
}
double distSegmentToSegment(Point a, Point b, Point c,
                            Point d) {
  if (segmentsIntersect(a, b, c, d)) return 0.0;
  return min({distPointToSegment(a, c, d),
              distPointToSegment(b, c, d),
              distPointToSegment(c, a, b),
              distPointToSegment(d, a, b)});
}
// rays a->b and c->d (not parallel)
bool raysIntersect(Point a, Point b, Point c, Point d) {
  Point ab = b - a, cd = d - c, ac = c - a;
  double det = cross(ab, cd);
  if (dcmp(det, 0) == 0) return false;
  double t = cross(ac, cd) / det;
  double u = cross(ac, ab) / det;
  return dcmp(t, 0) >= 0 && dcmp(u, 0) >= 0;
}
double distRayToRay(Point a, Point b, Point c, Point d) {
  if (raysIntersect(a, b, c, d)) return 0.0;
  return min(distPointToRay(a, c, d),
             distPointToRay(c, a, b));
}
// Line in the form a*x + b*y + c = 0
struct Line { double a, b, c; };
Line makeLine(Point p, Point q) {
  double a = p.y - q.y, b = q.x - p.x;
  return {a, b, -a * p.x - b * p.y};
}
double distToLine(Point p, Line l) {
  return fabs(l.a * p.x + l.b * p.y + l.c) /
         sqrt(l.a * l.a + l.b * l.b);
}
bool areParallel(Line l1, Line l2) {
  return fabs(l1.a * l2.b - l1.b * l2.a) < EPS;
}
// false if parallel, otherwise the point goes to res
bool intersectLines(Line l1, Line l2, Point &res) {
  double det = l1.a * l2.b - l2.a * l1.b;
  if (fabs(det) < EPS) return false;
  res.x = (l1.b * l2.c - l2.b * l1.c) / det;
  res.y = (l1.c * l2.a - l2.c * l1.a) / det;
  return true;
}
\end{lstlisting}
\subsection{Polygons}
\begin{lstlisting}
// Needs the two snippets above. Polygon = its corner
// points in order (CW or CCW).
double polygonArea(const vector<Point> &p) {
  double area = 0.0;
  int n = p.size();
  for (int i = 0; i < n; i++)
    area += cross(p[i], p[(i + 1) % n]); // > 0 if CCW
  return fabs(area) / 2.0;
}
// 1 = strictly inside, 0 = on the border, -1 = outside.
// Any simple polygon (ray casting). O(n)
int pointInPolygon(const vector<Point> &p, Point q) {
  int n = p.size();
  bool inside = false;
  for (int i = 0, j = n - 1; i < n; j = i++) {
    if (onSegment(p[j], q, p[i])) return 0;
    if ((p[i].y > q.y) != (p[j].y > q.y) &&
        q.x < (p[j].x - p[i].x) * (q.y - p[i].y) /
                (p[j].y - p[i].y) + p[i].x)
      inside = !inside;
  }
  return inside ? 1 : -1;
}
// Convex hull (monotone chain), CCW order. O(n log n).
// Collinear border points are removed; to KEEP them
// change both "<= 0" to "< 0".
vector<Point> convexHull(vector<Point> pts) {
  sort(all(pts));
  pts.erase(unique(all(pts)), pts.end());
  int n = pts.size(), k = 0;
  if (n <= 2) return pts;
  vector<Point> h(2 * n);
  for (int i = 0; i < n; ++i) { // lower hull
    while (k >= 2 && cross(h[k - 1] - h[k - 2],
                           pts[i] - h[k - 1]) <= 0)
      k--;
    h[k++] = pts[i];
  }
  for (int i = n - 2, t = k + 1; i >= 0; i--) { // upper
    while (k >= t && cross(h[k - 1] - h[k - 2],
                           pts[i] - h[k - 1]) <= 0)
      k--;
    h[k++] = pts[i];
  }
  h.resize(k - 1); // last point equals the first
  return h;
}
double perimeter(const vector<Point> &p) {
  double res = 0;
  int n = p.size();
  for (int i = 0; i < n; i++)
    res += dist(p[i], p[(i + 1) % n]);
  return res;
}
bool isConvex(const vector<Point> &p) {
  int n = p.size();
  if (n <= 2) return false;
  bool cw = false, ccw = false;
  for (int i = 0; i < n; i++) {
    int o = orientation(p[i], p[(i + 1) % n],
                        p[(i + 2) % n]);
    if (o == 1) cw = true;
    if (o == 2) ccw = true;
  }
  return !(cw && ccw);
}
Point polygonCentroid(const vector<Point> &p) {
  int n = p.size();
  Point c(0, 0);
  double area = 0.0;
  for (int i = 0; i < n; i++) {
    double cr = cross(p[i], p[(i + 1) % n]);
    area += cr;
    c = c + (p[i] + p[(i + 1) % n]) * cr;
  }
  return c / (3.0 * area);
}
// p = CONVEX polygon in CCW order (as from convexHull).
// 1 = inside, 0 = on the border, -1 = outside. O(log n)
int pointInConvexPolygon(const vector<Point> &p,
                         Point q) {
  int n = p.size();
  if (n < 3) return -1;
  if (cross(p[1] - p[0], q - p[0]) < -EPS ||
      cross(p[n - 1] - p[0], q - p[0]) > EPS)
    return -1;
  int l = 1, r = n - 1; // q is in wedge p[l], p[l+1]
  while (r - l > 1) {
    int mid = (l + r) / 2;
    if (cross(p[mid] - p[0], q - p[0]) >= 0) l = mid;
    else r = mid;
  }
  if (cross(p[r] - p[l], q - p[l]) < -EPS) return -1;
  if (onSegment(p[l], q, p[r])) return 0;
  if (l == 1 && onSegment(p[0], q, p[1])) return 0;
  if (r == n - 1 && onSegment(p[0], q, p[r])) return 0;
  return 1;
}
// Farthest pair of points (diameter). p = convex hull in
// CCW order. O(n) rotating calipers.
double diameter(const vector<Point> &p) {
  int n = p.size();
  if (n == 1) return 0;
  if (n == 2) return dist(p[0], p[1]);
  double best = 0;
  for (int i = 0, j = 1; i < n; i++) {
    Point e = p[(i + 1) % n] - p[i];
    // move j while it gets farther from the edge i, i+1
    while (cross(e, p[(j + 1) % n] - p[i]) >
           cross(e, p[j] - p[i]))
      j = (j + 1) % n;
    best = max({best, dist(p[i], p[j]),
                dist(p[(i + 1) % n], p[j])});
  }
  return best;
}
// is p strictly inside triangle abc?
bool pointInTriangle(Point p, Point a, Point b,
                     Point c) {
  int o1 = orientation(a, b, p);
  int o2 = orientation(b, c, p);
  int o3 = orientation(c, a, p);
  return o1 != 0 && o1 == o2 && o2 == o3;
}
// lattice (integer) points on segment ab, ends included
int latticeOnSegment(Point a, Point b) {
  int dx = abs((int)round(a.x - b.x));
  int dy = abs((int)round(a.y - b.y));
  return __gcd(dx, dy) + 1;
}
// Pick's theorem: Area = I + B/2 - 1, so the number of
// lattice points strictly inside is I = Area - B/2 + 1
// (polygon with integer vertices)
int latticeInside(const vector<Point> &p) {
  int n = p.size(), border = 0;
  for (int i = 0; i < n; i++)
    border += latticeOnSegment(p[i], p[(i + 1) % n]) - 1;
  return (int)round(polygonArea(p) - border / 2.0 + 1);
}
\end{lstlisting}
\subsection{Circles}
\begin{lstlisting}
// Needs "Point and vector basics" and "Lines".
struct Circle {
  Point center;
  double r;
};
// 0 = inside, 1 = on the circle, 2 = outside
int pointPosition(Circle c, Point p) {
  double d = dist(c.center, p);
  return dcmp(d, c.r) == 0 ? 1 : (d < c.r ? 0 : 2);
}
// circle through 3 points (not collinear): circumcircle
Circle circleFrom3Points(Point A, Point B, Point C) {
  double a2 = magSq(A), b2 = magSq(B), c2 = magSq(C);
  double D = 2 * (A.x * (B.y - C.y) + B.x * (C.y - A.y) +
                  C.x * (A.y - B.y));
  double ux = a2 * (B.y - C.y) + b2 * (C.y - A.y) +
              c2 * (A.y - B.y);
  double uy = a2 * (C.x - B.x) + b2 * (A.x - C.x) +
              c2 * (B.x - A.x);
  Point ctr(ux / D, uy / D);
  return {ctr, dist(ctr, A)};
}
// circle inside triangle abc touching all three sides
Circle incircle(Point a, Point b, Point c) {
  double sa = dist(b, c), sb = dist(a, c);
  double sc = dist(a, b);
  double per = sa + sb + sc;
  Point ctr = (a * sa + b * sb + c * sc) / per;
  double area = fabs(cross(b - a, c - a)) / 2.0;
  return {ctr, area / (per / 2.0)};
}
// centers of the circles of radius r through p1 and p2
vector<Point> circleCenters(Point p1, Point p2,
                            double r) {
  vector<Point> res;
  double d2 = magSq(p1 - p2), det = r * r - d2 / 4.0;
  if (det < 0.0 || d2 < EPS) return res;
  Point mid = (p1 + p2) / 2.0;
  Point shift = perp(p2 - p1) * (sqrt(det) / sqrt(d2));
  res.push_back(mid + shift);
  res.push_back(mid - shift);
  return res;
}
// LINE ab with circle c: 0, 1 or 2 points
vector<Point> lineCircle(Point a, Point b, Circle c) {
  vector<Point> res;
  Point p = projectOnLine(c.center, a, b);
  double h2 = magSq(p - c.center), r2 = c.r * c.r;
  if (dcmp(h2, r2) > 0) return res; // line misses
  double off = sqrt(max(0.0, r2 - h2));
  Point dir = normalize(b - a);
  res.push_back(p - dir * off);
  if (dcmp(off, 0) > 0) res.push_back(p + dir * off);
  return res;
}
vector<Point> segmentCircle(Point a, Point b, Circle c) {
  vector<Point> res;
  for (Point p : lineCircle(a, b, c))
    if (onSegment(a, p, b)) res.push_back(p);
  return res;
}
// intersection points of two circles: 0, 1 or 2 points
// (empty also when the circles are identical)
vector<Point> circleCircle(Circle c1, Circle c2) {
  vector<Point> res;
  double d = dist(c1.center, c2.center);
  if (dcmp(d, 0) == 0) return res; // same center
  if (dcmp(d, c1.r + c2.r) > 0 ||
      dcmp(d, fabs(c1.r - c2.r)) < 0)
    return res; // too far, or one inside the other
  // a = distance from c1 to the common chord
  double a =
    (c1.r * c1.r - c2.r * c2.r + d * d) / (2 * d);
  double h = sqrt(max(0.0, c1.r * c1.r - a * a));
  Point dir = (c2.center - c1.center) / d;
  Point mid = c1.center + dir * a;
  res.push_back(mid + perp(dir) * h);
  if (dcmp(h, 0) > 0) res.push_back(mid - perp(dir) * h);
  return res;
}
// touching points of the tangents from p to circle c
vector<Point> tangentsFromPoint(Point p, Circle c) {
  vector<Point> res;
  double d = dist(p, c.center);
  if (d < c.r - EPS) return res; // p is inside
  if (fabs(d - c.r) < EPS) return {p}; // p on the circle
  double alpha = acos(c.r / d); // angle at the center
  Point dir = (p - c.center) / d;
  res.push_back(c.center + rotate(dir, alpha) * c.r);
  res.push_back(c.center + rotate(dir, -alpha) * c.r);
  return res;
}
// Common tangents of two circles. Each pair = {touching
// point on c1, touching point on c2}. First the outer
// tangents, then the inner (crossing) ones.
vector<pair<Point, Point>> commonTangents(Circle c1,
                                          Circle c2) {
  vector<pair<Point, Point>> res;
  Point d = c2.center - c1.center;
  double dd = len(d);
  if (dd < EPS) return res; // same center
  Point dir = d / dd;
  if (dd >= fabs(c1.r - c2.r) - EPS) { // outer tangents
    double cs = (c1.r - c2.r) / dd;
    double al = acos(max(-1.0, min(1.0, cs)));
    for (double s : {al, -al}) {
      Point nrm = rotate(dir, s);
      res.push_back({c1.center + nrm * c1.r,
                     c2.center + nrm * c2.r});
    }
  }
  if (dd >= c1.r + c2.r - EPS) { // inner tangents
    double al = acos(min(1.0, (c1.r + c2.r) / dd));
    for (double s : {al, -al}) {
      Point nrm = rotate(dir, s);
      res.push_back({c1.center + nrm * c1.r,
                     c2.center - nrm * c2.r});
    }
  }
  return res;
}
// area of the intersection of two circles
double circleIntersectionArea(Circle c1, Circle c2) {
  double d = dist(c1.center, c2.center);
  if (d >= c1.r + c2.r) return 0.0;
  if (d <= fabs(c1.r - c2.r)) { // one inside the other
    double r = min(c1.r, c2.r);
    return PI * r * r;
  }
  double r1 = c1.r, r2 = c2.r;
  // central angles of the two circular segments
  double a1 =
    2 * acos((r1 * r1 + d * d - r2 * r2) / (2 * r1 * d));
  double a2 =
    2 * acos((r2 * r2 + d * d - r1 * r1) / (2 * r2 * d));
  return 0.5 * r1 * r1 * (a1 - sin(a1)) +
         0.5 * r2 * r2 * (a2 - sin(a2));
}
// theta = central angle in radians:
//  chord length   = 2 * r * sin(theta / 2)
//  angle of chord = 2 * asin(chord / (2 * r))
//  arc length     = r * theta
//  sector area    = r * r * theta / 2
//  segment area   = r * r * (theta - sin(theta)) / 2
\end{lstlisting}
\subsection{Integer geometry (exact)}
\begin{lstlisting}
// Use this (no doubles, no EPS) when coordinates are
// integers. Coordinates up to ~1e9 fit in ll.
struct P {
  ll x, y;
  P operator-(const P &o) const {
    return {x - o.x, y - o.y};
  }
};
ll cross(P a, P b) { return a.x * b.y - a.y * b.x; }
ll dot(P a, P b) { return a.x * b.x + a.y * b.y; }
ll distSq(P a) { return a.x * a.x + a.y * a.y; }
// > 0: a->b->c turns left (CCW), < 0: right, 0: line
ll ccw(P a, P b, P c) { return cross(b - a, c - a); }
// Sort by angle around the origin, from 0 to 360 degrees
// (starting at the positive x-axis), ties by distance.
// Usage: sort(all(v), angleCmp);  (around point c: use
// a - c and b - c)
bool upper(P p) {
  return p.y > 0 || (p.y == 0 && p.x > 0);
}
bool angleCmp(P a, P b) {
  if (upper(a) != upper(b)) return upper(a);
  ll cp = cross(a, b);
  if (cp != 0) return cp > 0;
  return distSq(a) < distSq(b);
}
// same order as atan2: from -180 (exclusive) to +180
bool lower(P p) {
  return p.y < 0 || (p.y == 0 && p.x > 0);
}
bool atan2Cmp(P a, P b) {
  if (lower(a) != lower(b)) return lower(a);
  ll cp = cross(a, b);
  if (cp != 0) return cp > 0;
  return distSq(a) < distSq(b);
}
// Sort points CCW around point L, starting from the
// direction of the ray L -> R. Set L and R first.
P L, R;
int region(P p) { // 0: on the ray, 1: left, 2: behind, 3
  ll c = ccw(L, R, p);
  if (c > 0) return 1;
  if (c < 0) return 3;
  return dot(R - L, p - L) >= 0 ? 0 : 2;
}
bool polarCmp(P a, P b) {
  int ra = region(a), rb = region(b);
  if (ra != rb) return ra < rb;
  ll c = ccw(L, a, b);
  if (c != 0) return c > 0;
  return distSq(a - L) < distSq(b - L);
}
// reduce a direction to its smallest integer form, e.g.
// (4, -6) -> (2, -3). Good as a map key for slopes.
P simplify(P p) {
  ll g = __gcd(abs(p.x), abs(p.y));
  if (g == 0) return p;
  p.x /= g, p.y /= g;
  if (p.x < 0 || (p.x == 0 && p.y < 0))
    p.x = -p.x, p.y = -p.y;
  return p;
}
\end{lstlisting}
\subsection{Closest pair and enclosing circle}
\begin{lstlisting}
// Needs "Point and vector basics" and "Circles".
// Smallest distance between two points. O(n log n)
double closestPair(vector<Point> p) {
  sort(all(p),
       [](Point a, Point b) { return a.x < b.x; });
  multiset<pair<double, double>> box; // {y, x} of nearby
  double best = 1e18;
  for (int i = 0, j = 0; i < (int)p.size(); i++) {
    // drop points too far to the left
    while (j < i && p[i].x - p[j].x >= best) {
      box.erase(box.find({p[j].y, p[j].x}));
      j++;
    }
    auto it = box.lower_bound({p[i].y - best, -1e18});
    for (; it != box.end() && it->first <= p[i].y + best;
         ++it)
      best =
        min(best, dist(p[i], {it->second, it->first}));
    box.insert({p[i].y, p[i].x});
  }
  return best;
}
// Smallest circle that contains all points (Welzl).
// Expected O(n) after the shuffle. Uses rng.
Circle minEnclosingCircle(vector<Point> p) {
  shuffle(all(p), rng);
  int n = p.size();
  Circle c{p[0], 0};
  auto out = [&](Point q) {
    return dist(c.center, q) > c.r + EPS;
  };
  for (int i = 1; i < n; i++) {
    if (!out(p[i])) continue;
    c = {p[i], 0}; // p[i] must be on the border
    for (int j = 0; j < i; j++) {
      if (!out(p[j])) continue;
      c = {(p[i] + p[j]) / 2, dist(p[i], p[j]) / 2};
      for (int k = 0; k < j; k++)
        if (out(p[k]))
          c = circleFrom3Points(p[i], p[j], p[k]);
    }
  }
  return c;
}
\end{lstlisting}
\end{multicols*}
\end{document}
